\documentclass[a4paper,12pt]{book}%{report}
\usepackage{microtype}
%
\usepackage{a4wide}
\usepackage[amssymb]{SIunits}
\usepackage{units}
\usepackage{pslatex}
\usepackage{amssymb}
\usepackage{amsbsy}
\usepackage[fleqn]{amsmath}
\usepackage{bbm}
\usepackage{mathtools}
\usepackage{booktabs}
\usepackage{url}
\usepackage[many]{tcolorbox} %Achtung! Reihenfolge beachten
\usepackage{xcolor}
\usepackage[utf8]{inputenc}
\usepackage[ngerman]{babel} % needs: sudo apt-get install texlive-lang-german
\usepackage[figurename=Abb.]{caption}
\usepackage{calrsfs}
\usepackage[titletoc]{appendix}
\usepackage{thumbpdf}
\usepackage[colorlinks,
pdfpagelabels,
pdfstartview = FitH,
bookmarksopen = true,
bookmarksnumbered = true,
linkcolor = black,
plainpages = false,
hypertexnames = false,
citecolor = black]{hyperref}
% Define the headers
\usepackage{fancyhdr}
\pagestyle{fancy}
\fancyhf{}
\fancyhead[EL]{\leftmark}
\fancyhead[OR]{\rightmark}
\fancyfoot[OR]{\thepage}
\fancyfoot[EL]{\thepage}
\renewcommand{\headrulewidth}{0pt}
% This is to have text flow around figures
\usepackage{wrapfig}
% Use tikz and pgfplot thru gnuplot
\usepackage[miktex]{gnuplottex}
\usepackage{gnuplot-lua-tikz}
\usepackage{tikz}
\usepackage{pgfplots}
%\pgfplotsset{compat=1.14}
\usetikzlibrary{decorations.pathreplacing,decorations.markings} %,snakes}
\usetikzlibrary{shapes,arrows,chains,matrix,positioning,scopes}
\usetikzlibrary{positioning}
\usetikzlibrary{calc}
\usetikzlibrary{arrows.meta}
\tikzset{darkstyle/.style={circle,draw,fill=gray!40}}
%colours
\definecolor{darkgreen}{rgb}{0,0.6,0}
\newcommand{\dgreen}{\color{darkgreen}}
\newcommand{\red}{\color{red}}
\newcommand{\blue}{\color{blue}}
\newcommand{\cyan}{\color{cyan}}
\newcommand{\green}{\color{green}}
\newcommand{\mbs}[1]{\ensuremath{\boldsymbol{\mathbf{#1}}}}
% New command and renewcommand settings
\newcommand{\mbf}{\mathbf}
\newcommand{\bs}{\boldsymbol}
%\renewcommand{\thepage}{{\thepart}.\arabic{page}}
%\addto\captionsgerman{\renewcommand{\figurename}{Abb.}}
\newtheorem{theorem}{Satz}
% New environment settings
%
\definecolor{myblue}{RGB}{0,163,243}
\tcbset{mystyle/.style={ 
%  breakable, 
  enhanced, 
  outer arc=0pt, 
  arc=0pt, 
  colframe=gray, 
  colback=white, 
  boxed title style={ 
    colback=gray,
    coltitle=black,
    colbacktitle=gray,
  }, 
  outer arc=5pt,
  title=Boxhead~\thetcbcounter, 
  fonttitle=\sffamily 
  } 
}
%
\newtcolorbox[auto counter]{satz}[1]{mystyle, title=Satz~\thetcbcounter: #1}
\newtcolorbox[auto counter]{exercise}[1]{mystyle, title=Aufgabe~\thetcbcounter: #1}
\newtcolorbox[auto counter]{example}[1]{breakable,mystyle, title=Beispiel~\thetcbcounter: #1}
\newtcolorbox{note}[1]{ enhanced,attach boxed title to top left={yshift=-3mm,yshifttext=-1mm},
  title={N.B.: #1},fonttitle=\bfseries,
  boxed title style={size=small}}
%
\newboolean{showComments} %Textblöcke anzeigen
\newcommand\Comment[1]{\ifthenelse{\boolean{showComments}}{\leavevmode\newline{\blue #1}\newline}{}}
%
\newboolean{showErratum} %Textblöcke anzeigen
\newcommand\Erratum[1]{\ifthenelse{\boolean{showErratum}}{\leavevmode\newline{\red #1}\newline}{}} 
%
%\title{\huge Lecture Notes\\
%\large für Studierende der Studiengänge Embedded Systems %Engineering, Mirkosystemtechnik und Sustainable Systems %Engineering}
%\author{Andreas Greiner und Lars Pastewka \\Professur für %Simulation\\Institut für Mikrosystemtechnik\\Universität %Freiburg}
\parindent0em
\parskip1ex
\pagestyle{headings}
\setcounter{secnumdepth}{3}
% The following is better placed in the main document
%\setboolean{showComments}{true}
%\setboolean{showErratum}{true}
\parindent0em
\parskip1ex
\pagestyle{headings}
\setboolean{showComments}{false}
\setboolean{showErratum}{true}
%
\title{\huge Lecture Notes\\
\large für Studierende der Studiengänge Embedded Systems Engineering, Mirkosystemtechnik und Sustainable Systems Engineering}
\author{Andreas Greiner und Lars Pastewka \\Professur für Simulation\\Institut für Mikrosystemtechnik\\Universität Freiburg\\[1cm]
  \includegraphics[width=0.4\textwidth]{fig/Uni_Logo-Grundversion_E1_A4_RGB}
}
%
%\usepackage{caption}
%\usepackage{lipsum}
%\usepackage{cutwin}
%
\begin{document}
\frontmatter
\maketitle
\part{Differentialgleichungen}
%\Erratum{Hier ist eine Fehlerbeschreibung.}
Lernziele:
\begin{itemize}
\item Verständnis der Klassifizierung unterschiedlicher Differenzialgleichungen
\item Benutzung unterschiedlicher Typen von Gleichungen zur Modellierung von Phä\-nomenen
\item Analytische und numerische Lösungen und deren Visualisierung und Interpretationen
\end{itemize}
Die meisten Phänomene denen wir in den Ingenieurwissenschaften begegnen, werden
sehr gut durch Differentialgleichungen beschrieben. Lineare zeitinvariante
System der Systemtheorie {(LTI-Systeme)} im Zeitbereich, 
%\Comment{ anderes Beispiel wählen!!!} 
werden durch ein System linearer gewöhnlicher
Differentialgleichungen (ODE = ordinary differential equation) mit der Zeit als
unabängiger Veränderlicher beschrieben. Ein Diffusionsprozess, wie z.B. der
Wärmetransport in einem Bauteil auf einem Kühlkörper, das einer Wärmequelle
ausgesetzt ist, ist ein Phänomen das am Besten mit einer partiellen
Differentialgleichung (PDE = partial differential equation) beschrieben wird.

Literatur:
\begin{itemize}
  \item J. Brenner und P. Lesky, Mathematik für Ingenieure und
    naturwissenschaftler I-IV. Dies ist ein sehr umfassendes Werk mit vielen
    Beispielen. Band III enthält Kapitel über Differentialgleichungen, Band IV
    übe Fourier- und Laplacetransformationen aber auch vieles mehr. Einer der
    Autoren dieses Manuskripts (A.G.) hatte die Ehre die Vorlesungsreihe bei
    Peter Albin Lesky\footnote{Peter Albin Lesky (* 6. Dezember 1926 in Graz; †
    12. Februar 2008 in Innsbruck) war ein österreichischer Mathematiker
    (https://de.wikipedia.org/wiki/Peter\_Albin\_Lesky). } hören zu dürfen und ist
    heute noch von seiner didaktischen Kunst beeindruckt.
  \item V. I. Arnol'd, Gewöhnliche Differentialgleichungen. Formal, kurze
    Darstellung.
  \item G. Baker, Differential equations as models in science and engineering.
    Sehr viele anwendungsbezogene Beispiele.
  \item Frank Ayres, Differentialgleichungen. 1900 ausführliche Lösungsbeispiele.
  \item J. Jäckle, Einführung in die Transporttheorie. Dies wurde als Grundlage
    für das Kapitel über die Theorie der linearen Antwort genommen.
  \item H. Haken, Advanced Synergetics, 2nd printing, Springer 1987. Für das
    Kapitel über die Adiabatische Elimination in Systemen nichtlinearer
    Differentialgleichungen. Diese Buch geht weit über das hier dargestellte
    hinaus und behandelt Anwendungen des von H. Haken entwickelten Formalismus
    auf zahlreiche naturwissenschaftliche Phänomene. Besonders zu empfehlen.
  \item Mary L. Boas,  Mathematical Methods in the Physical Sciences, Wiley
    2005. Sehr anwendungsorientiert mit vielen Anwendungsbeispielen.
    Kompendium zu Höherer Mathematik 
  \item K. Meyberg und P. Vachenauer, Höhere Mathematik 2, Springer 1997.  Sehr
    mathematisch gehalten mit Merkboxen zum Vorgehen bei der Lösung von
    Differentialgleichungen. 
  \item Harumi Hattori, Partial differential equations: methods, applications
    and theories, World Scientific 2019
\end{itemize}
Dies ist eine sehr eingeschränkte Sicht auf die Literatur. Diese Liste wird,
wie das ganze Manuskript, ständig erweitert werden.
\newpage
\tableofcontents
\newpage 
\mainmatter
\chapter[Erinnerung an die \dots]{Erinnerung an die Integral- und
Differentialrechnung und die Lineare Algebra}
In diesem Kapitel wollen wir an die in den Grundvorlesungen der Mathematik für
Ingenieurinnen und Informatiker behandelten Themen erinnern.
\section{Sätze der Differential- und Integralrechnung}
\begin{satz}{Hauptsatz der Differential- und Integralrechnung\label{theo:HSIntegralDiff}}
Wenn für eine stetige Funktion $f(x)$ das Integral
$F(x)=\int\limits_{x_0}^{x}f(\xi)d\xi$ existiert, dann ist
\[
  \frac{dF(x)}{dx}=\frac{d}{dx}\int\limits_{x_0}^{x}f(\xi)d\xi=f(x)
\]
\label{theo:HDI}
\end{satz}
Satz \ref{theo:HDI} in Worten:„Die Ableitung des Integrals ist gleich der
Funktion im Integranden an der oberen Integrationsgrenze“.
\begin{example}{Ableitung eines Integrals}
\[
 \frac{d F(x)}{dx}=\frac{d}{dx}\int\limits_{x_0}^{x}e^{-\xi}d\xi=
 \frac{d}{dx} \left(-e^{-x}+e^{-x_0}\right)=e^{-x}
\]
\end{example}
\begin{satz}{Mittelwertsatz der Integralrechnung}
  Gegeben eine Funktion $f(x)$, die auf dem Interval $[a,b]$ stetig ist, sowie
  eine Funktion $g(x)$ für die im Intervall $[a,b]$ entweder $g(x)\ge0$ oder
  $g(x)\le0$ gelte. Dann gibt es mindestens eine Stelle $\xi\in[a,b]$, so dass
  mit deren entsprechendem Funktionswert 
  \begin{equation} 
  \int\limits_a^bf(x)g(x)dx=f(\xi)\int\limits_a^bg(x)dx 
  \label{eq:MittelwertInt2} 
  \end{equation}
  gilt.   
\end{satz}
Wenn wir $g(x)=1 $ setzen, dann wird dies gemeinhin als erster
Mittelwertsatz der Integralrechnung bezeichnet. Es gilt also:
\begin{equation} 
  \int\limits_a^bf(x)dx=f(\xi)(b-a)
  \label{eq:MittelwertInt1}
\end{equation}
%\begin{example}{Anwendung von Satz \ref{eq:MittelwertInt2}}
%?
%\end{example}
\begin{satz}{Mittelwertsatz der Differentialrechnung}
  Gegeben 2 Funktionen $f(x)$ und $g(x)$, die beide auf dem Interval $[a,b]$
  stetig und differenzierbar seien. Dann existiert ein $x_0\in(a,b)$, so dass   
  \begin{equation}
    f'(x_0)\left(g(b)-g(a)\right)= g'(x_0)\left(f(b)-f(a)\right)
  \label{eq:MittelwertDiff2}
\end{equation}
\end{satz}
Wenn wir $g'(x)=1 $ setzen, dann wird dies gemeinhin als erster
Mittelwertsatz der Differetialrechnung bezeichnet. Es gilt also:
\begin{equation}
  f'(x_0)=\frac{f(b)-f(a)}{b-a}
  \label{eq:MittelwertDiff1}
\end{equation}
\section{Taylorreihe}
%{\red Das könnte ein Beispiel für den Beweis eines Satzes sein}
Die Taylorreihe ist zweifelsohne eines der wichtigen Instrumente. Sie dient zur
Abschätzung des Verhaltens von Funktionen und findet in numerischen
Methoden Anwendung, wie z.B. den Finiten Differenzen.
\begin{satz}{Der Taylorsche Satz\label{theo:Taylor}} Sei $f$ eine Funktion, die
  eine $(n+1)$-te stetige Ableitung auf einem Intervall $J$ besitze. Es seien
  $a,b\in J$. Dann ist 
\[
    f(b)=f(a)+\frac{b-a}{1!}f'(a)+\cdots +\frac{(b-a)^{n}}{(n)!}f^{(n)}(a)+R_n
\]
  mit
\[
    R_n=\int\limits_a^b\frac{(b-s)^{n}}{n!}f^{(n+1)}(s)ds
\]
Des weiteren merken wir an, dass eine Zahl $c\in[a,b]$ existiert, so dass
\[
  R_n=\frac{(b-a)^{n+1}}{(n+1)!}f^{(n+1)}(c)
\]
gilt.
\end{satz}
{\bf Beweis:} Wir wenden Satz 1 an und integrieren partiell
\begin{align}
  f(b)=&f(a)+\int\limits_a^bf'(s)ds=f(a)+\int\limits_a^b\frac{(b-s)^0}{0!}f'(s)ds\nonumber\\
  =&f(a)-\left.\frac{(b-s)^1}{1!}f'(s)\right|_a^b+\int\limits_a^b\frac{(b-s)^1}{1!}f''(s)ds=\dots
  \label{eq:TaylorProof1}
\end{align}
Vollständige Induktion liefert den Schluss des Beweises.
\begin{example}{Taylorreihenentwicklung}
Man entwickle $f(x)=e^{-x}$ um $x=0$
  \[ e^{-x}=1-x+\frac{1}{2!}x^2-\frac{1}{3!}x^3+R_3 \]
\end{example}
\section{Lineare Algebra}
In der Mathematik treffen wir auf verschiedene Arten von Objekten, die
untereinander addiert und mit Zahlen multipliziert werden können. Eine
Ansammlung von Objekten nennen wir eine Menge. Ein Beispiel ist die Menge von
Vektoren ein und derselben Dimension. Wir wollen für all die
Spezialfälle eine allgemeine Definition geben. 
\subsection{Der Vektorraum}
Ein {\bf Vektorraum} V ist eine Menge von Objekten, die addiert und mit Zahlen
multipliziert werden können und zwar so, dass die Summe zweier Elemente aus V
wieder ein Element aus V ergibt, das Produkt eines Elementes aus V mit einer Zahl
ebenfalls wieder in V ist und folgende Eigenschaften erfüllt sind:.
\begin{description}
  \item[{\bf VR 1.}] Gegeben $u,v,w\in V$. Es soll Assoziativität gelten: $(u+v)+w = u+(v+w)$
  \item[{\bf VR 2.}] Es gibt ein Element 0 in V, so dass gilt $0+u = u+0 =u$
    für alle $u\in V$. Wir nennen dieses Element den Nullvektor.
  \item[{\bf VR 3.}] Für jedes $u\in V$ gibt es ein Element $(-1)u\in V$ für das $u+(-1)u=0$
    gilt. D.h. jedes Element hat ein Inverses.
  \item[{\bf VR 4.}] Für alle $u,v\in V$ soll Kommuntativität gelten $u+v = v+u$
  \item[{\bf VR 5.}] Für eine Zahl $c$ gilt Distributivität $c(u+v)=cu+cv$ 
  \item[{\bf VR 6.}] Für zwei Zahlen $a$, $b$ gilt Assoziativität bezüglich der
    Addition $(a+b)u=au+bu$
  \item[{\bf VR 7.}] Für zwei Zahlen $a$, $b$ gilt Assoziativität bezüglich der
    Multiplikation $(ab)u=a(bu)$
  \item[{\bf VR 8.}] Für alle $u\in V$ gilt $1\cdot u=u$
\end{description}
Als einen Unterraum $W$ von $V$ bezeichnen wir einen Raum für den gilt
\begin{enumerate}
  \item $\forall u,v\in W$ ist auch $u+v\in W$
  \item $\forall u\in W$ und $c\in\mathbb{R}$ ist auch $cu\in W$
  \item $0\in V$ ist auch Element von $W$ 
\end{enumerate}
\begin{example}{Beispiele für Vektorräume}
  \begin{itemize}
    \item Der dreidimensionale Euklidische Raum $\mathbf{R}^3$ ist ein Vektorraum.
    \item Die Menge der Polynome kleiner gleich vierten Grades 
      \[p_4(x)=a_0+a_1x+a_2x^2+a_3x^3+a_4x^4\]
      ist ein Vektorraum.
  \end{itemize}
  Man verifiziere dies anhand der oben gegebenen Definition.
\end{example}
 \subsection{Lineare Unabhängigkeit, Basis}
Eine Linearkombination von Vektoren $(v_1,\dots,v_m)\in V$ schreiben wir als
$a_1v_1+\dots+a_mv_m$, wobei $a_1,\dots,a_m\in\mathbb{R}$ oder auch
$a_1,\dots,a_m\in\mathbb{C}$. Alle möglichen Linearkombinationen der $v_i$
nennen wir die lineare Hülle der $(v_1,\dots,v_m)\in V$ und schreiben dafür
$span(v_1,\dots,v_m)$. Die $m$ Vektoren spannen einen Unteraum $W\subset V$
auf.
\begin{example}{Lineare Hülle}
Es ist
\[(7,2,9)=2(2,1,3)+3(1,0,1). \]
Der Vektor $(7,2,9)$ ist also eine Linearkombination von $(2,1,3)$ und
$(1,0,1)$.  Daher sagen wir auch $(7,2,9)\in span\left( (2,1,3),(1,0,1)
\right)$ und damit nennen wir diesen Vektor auch linear abhängig.
\end{example}
Wir nennen Vektoren $(v_1,\dots,v_m)\in V$ linear unabhängig, wenn keiner der
$m$ Vektoren als Linearkombination der restlichen $m-1$ Vektoren geschrieben
werden kann. Eine Basis des Vektorraums $V$ ist gegeben durch die maximale
Anzahl linear unabhängiger Vektoren, die ganz $V$ aufspannen. Diese muss nicht
eindeutig sein.
\begin{example}{Basen}
  \begin{itemize}
    \item $(1,0,0)$, $(0,1,0)$ und $(0,0,1)$ spannen der $\mathbb{R}^3$ auf. 
    \item Genauso tut das aber auch die Basis
      $(1/\sqrt{2},-1/\sqrt{2},0)$,$(-1/\sqrt{2},1/\sqrt{2},0)$ und $(0,0,1)$.
    \item Die Monome $\{1,x,x^2,x^3,x^4\}$ bilden die Basis der Vektorraums aus
      der Menge der Polynome vierten Grades und kleiner.
  \end{itemize}
\end{example}
\subsection{Der Rang einer Matrix}
Gegeben die Matrix
\begin{equation}
  \mathbf{A}=
  \begin{pmatrix}
    a_{11}&a_{12}&\;\cdots\;&a_{1n}\\
    a_{21}&a_{22}&\;\cdots\;&a_{2n}\\
    & &\;\cdots\;&\\
    a_{m1}&a_{m2}&\;\cdots\;&a_{mn}\\
  \end{pmatrix}
  \label{eq:Matrixmn}
\end{equation}
\begin{note}{}
  Vektoren und Matrizen schreiben wir in Fettdruck. Wir bezeichnen der Rang
  einer Matrix $\mathbf{A}$ mit $\text{Rg}(\mathbf{A})$.
\end{note}
Durch Streichen von Zeilen oder Spalten erhalten wir quadratische $s\times s$
Untermatrizen. Gehen wir davon aus, dass $A$ nicht die Nullmatrix ist, dann finden
wir sicherlich unter den $s\times s$ Untermatrizen solche, deren Determinanten
von Null verschieden sind. Die Maximale Zahl $r$ an Reihen, bzw.  Spalten, der
von Null verschiedenen Determinanten, die bei dieser Operation entstehen,
nennen wir den Rang der Matrix $\text{Rg}(\mathbf{A})$.
\begin{example}{Rang einer $5\times7$-Matrix}
  Man bestimme den Rang der Matrix
  \[
    A=\begin{pmatrix}
    7&1&0&2&-1&4&5\\
    1&1&2&3& 0&1&2\\
    0&1&-2&1&2&0&1\\
    4&-1&-8&-6&1&1&0\\
    0&1&2&1&4&0&1\\
    \end{pmatrix}
  \]
  Wir wissen aus der linearen Algebra, dass der Wert einer Detrminante
  unverändert bleibt, wenn zu einer Zeile bzw. einer Spalte ein beliebiges
  Vielfaches einer Zeile bzw. einer Spalte addiert wird.

  Dies machen wir uns zunutze und addieren Vielfache der Zeilen der obigen
  Matrix sukzessive solange, bis folgende Matrix entsteht: 
  \[
    \begin{pmatrix}
    1&0&0&0&0&0\\
    0&1&0&0&0&0\\
    0&0&1&-2&-2&1\\
    0&0&0&13&11&-3\\
    \end{pmatrix}
  \]
  Beachte: alle linear abhängigen Zeilen und Spalten können wir streichen. 
  
  Als letzte Umformung erhalten wir
  \[
     \begin{pmatrix}
    1&0&0&0\\
    0&1&0&0\\
    0&0&1&0\\
    0&0&0&1\\
    \end{pmatrix}
  \]
  woraus wir den Rang $r=\text{Rg}(\mathbf{A})=4$ ablesen.
\end{example}
Der Rang einer Matrix $\mathbf{M}$ ist die maximale Anzahl linear unabhängiger
Zeilen bzw. Spalten. Die linear unabhängigen Zeilenvektoren spannen die lineare
Hülle des Zeilenvektorraums auf.
%
\subsection{Lineare Gleichungssysteme (LGS)}
Nach dem Satz von Kronecker-Capelli ist ein LGS 
$\mathbf{A\cdot x}=\mathbf{b}$ genau dann lösbar, wenn der Rang der
Koeffizientenmatrix $\mathbf{A}$ gleich dem Rang der erweiterten
Koeffizientenmatrix $\mathbf{A|b}$ ist. Die erweiterte Koeffizientenmatrix
$\mathbf{A|b}$ bilden wir, indem wir den Spaltenvektor $\mathbf{b}$ an die
Matrix $\mathbf{A}$ anfügen. Wir schreiben dies $\mathbf{\{A|b\}}$.
\begin{example}{Lösbarkeit eines LGS}
  \begin{itemize}
    \item
      \[
	\underbrace{\begin{pmatrix}1&1\\2&2\end{pmatrix}}_{\mathbf{A}}
        \underbrace{\begin{pmatrix}x_1\\x_2\end{pmatrix}}_{\mathbf{x}}=
        \underbrace{\begin{pmatrix}2\\4\end{pmatrix}}_{\mathbf{b}}
	\qquad\text{ also schreiben wir }\mathbf{A\cdot x}=\mathbf{b}
      \]
      Es ist $\text{Rg}(\mathbf{A})=1$ und $\text{Rg}(\mathbf{\{A|b\}})=1$,
      also ist das LGS lösbar. Allerdings erhalten wir in diesem Fall eine
      paramterische Lösung und es existieren somit unendlich viele Lösungen.
    \item
       \[
        \underbrace{\begin{pmatrix}1&1\\2&2\end{pmatrix}}_{\mathbf{A}}
        \underbrace{\begin{pmatrix}x_1\\x_2\end{pmatrix}}_{\mathbf{x}}=
        \underbrace{\begin{pmatrix}2\\3\end{pmatrix}}_{\mathbf{c}}
        \qquad\text{ also schreiben wir }\mathbf{A\cdot x}=\mathbf{c}
      \]
      Es ist $\text{Rg}(\mathbf{A})=1$ und $\text{Rg}(\mathbf{\{A|c\}})=2$,
      also ist das LGS nicht lösbar. In diesem Fall sehen wir das auch gleich
      mit bloßem Augen, denn die beiden Gleichungen $x_1+x_2=2$ und
      $x_1+x_2=1.5$ führen auf einen Widerspruch.

      Wohingegen das LGS
    \item
       \[
        \underbrace{\begin{pmatrix}1&1\\1&2\end{pmatrix}}_{\mathbf{D}}
        \underbrace{\begin{pmatrix}x_1\\x_2\end{pmatrix}}_{\mathbf{x}}=
        \underbrace{\begin{pmatrix}2\\3\end{pmatrix}}_{\mathbf{c}}
	\qquad(\text{wir schreiben }\mathbf{D\cdot x}=\mathbf{c})
      \]
      eine eindeutige Lösung hat. Der Nachweis sei dem geneigten Leser überlassen.

      \textbf{Anleitung}: Berechne zunächst zur Übung $\text{Rg}(\mathbf{D})$ und
      $\text{Rg}(\mathbf{\{D|c\}})$. Berechne dann die Lösung.


  \end{itemize}
\end{example}
\subsection{Eigenwerte, Eigenvektoren, Normalformen}
Die Menge der Eigenwerte und Eigenvektoren einer Matrix nennt man auch ihr Spektrum. Von den drei Matrizen $\mathbf{M}_i$ mit $i=1,2,3$
\[
  \mathbf{M}_1=\begin{pmatrix}
    -9&2&6\\5&0&-3\\-16&4&11
  \end{pmatrix},\quad
  \mathbf{M}_2=\begin{pmatrix}
    2&4&-2\\4&2&-2\\-2&-2&-1
  \end{pmatrix}\text{ und }
  \mathbf{M}_3=\begin{pmatrix}
    6&-6&5\\14&-13&10\\7&-6&4
  \end{pmatrix}
\]
wollen wir das Spektrum berechnen. Es ist $\text{Rg}(\mathbf{M}_i)=3$ (Nachweis!).
\[
  \begin{vmatrix}
    -9-\lambda&2&6\\5&0-\lambda&-3\\-16&4&11-\lambda
  \end{vmatrix}
  =-(\lambda-1)\cdot(\lambda+1)\cdot(\lambda-2)=0
\]
Wir erhalten drei verschiedene Eigenwerte $\lambda_1=1$, $\lambda_2=-1$ und $\lambda_3=2$. Aus $(\mathbf{M}_1-\lambda_i\mathds{1})\cdot\mathbf{x}_i=0$ berechnen wir die Eigenvektoren.  Für $\lambda_1=1$ erhalten wir das LGS \[
 \begin{pmatrix}-10&2&6\\5&-1&-3\\-16&4&10\end{pmatrix}\cdot
 \begin{pmatrix}x_{11}\\x_{12}\\x_{12}\end{pmatrix}=
 \begin{pmatrix}0\\0\\0\end{pmatrix}
\]
Dies ist ein LGS, dessen Rang der Koeffizientenmatrix 2 ist, gleich dem Rang der erweiterten Koeffizientenmatrix. Damit erwarten wir eine parametrische Lösung mit einem frei wählbaren Parameter.
Wir sehen, dass die ersten beiden Zeilen der Koeffizientenmatrix linear abhängig sind (multipliziere 2. Zeile mit $-2$ und erhalte Zeile 1). Davon ist allerdings die dritte Zeile unabhängig. Also hat die Koeffizientenmatrix den Rang 2. Der Rangabfall gegenüber $\mathbf{M}_1$ ist 1 und dies ist gleich der Vielfachheit des Eigenwertes, im voeliegenden Fall 1. In der Tat haben wir 3 verscheidene Eigenwerte.

Die Eigenvektoren $\mathbf{x}_i$ lauten daher
\[
\mathbf{x}_1=\begin{pmatrix}1\\ -1\\ 2\end{pmatrix},\quad
\mathbf{x}_2=\begin{pmatrix}2\\ -1\\ 3\end{pmatrix}\text{ und }
\mathbf{x}_3=\begin{pmatrix}2\\ -1\\ 4\end{pmatrix}.
\]
Wir haben die Komponenten so gewählt, dass sie ganze Zahlen ergeben, das ist natürlich will\-kürlich bei einer Parameterlösung.

Die Eigenvektoren bilden die Transformationsmatrix $\mathbf{P}_1$, mit deren Hilfe wir $\mathbf{M}_1$ auf Diagonalgestalt bringen und zwar mit $\bm{\Lambda}_1=
\mathbf{P}_1^{-1}\mathbf{M}_1\mathbf{P}_1$. Es ist
\[
\mathbf{P}_1=\begin{pmatrix}1&2&2\\ -1&-1&-1\\ 2&3&4\end{pmatrix}
\quad\text{ und }\quad
\mathbf{P}_1^{-1}=\begin{pmatrix}-1&-2&0\\ 2&0&-1\\ -1&2&1\end{pmatrix}.
\]
Wir erhalten somit
\[
\bm{\Lambda}_1=
\begin{pmatrix}-1&-2&0\\ 2&0&-1\\ -1&2&1\end{pmatrix}
\begin{pmatrix}-9&2&6\\5&0&-3\\-16&4&11\end{pmatrix}
\begin{pmatrix}1&2&2\\ -1&-1&-1\\ 2&3&4\end{pmatrix}=
\begin{pmatrix}1&0&0\\0&-1&0\\0&0&2\end{pmatrix}
\]

Die Eigenwerte von $\mathbf{M}_2$ sind $\lambda_1=7$ und $\lambda_2=\lambda_3=-2$. wir haben einen doppelt auftretenden Eigenwert. Für $\lambda_1=7$ erhalten wir den Eigenvektor $\mathbf{x}_1=(2/3,2/3,-1/3)^T$. Was passiert nun mit dem doppelten Eigenwert, können wir hierfür Eigenvektoren identifizieren? Wir untersuchen die Matrix  $\mathbf{M}_2-\lambda_{2}\mathds{1}$ und stellen fest dass diese einen Rangabfall von 2 aufweist. Damit erhalten wir für die Lösung des Gleichungssystems $(\mathbf{M}_2-\lambda_{2}\mathds{1})\cdot\mathbf{x}_i=0$ eine zweiparametrische Lösung und wir finden die zwei linear unabhängigen Eigenvektoren
\[
\mathbf{x}_2=\frac{\sqrt{2}}{2}
\begin{pmatrix}1\\ 1\\ 0\end{pmatrix}
\quad\text{ und }
\mathbf{x}_3=\frac{\sqrt{2}}{6}
\begin{pmatrix}1\\ 1\\ 4\end{pmatrix}
\]
Die Konstruktion von $\mathbf{P}_2$ und $\mathbf{P}_2^{-1}$ mit diesen Eigenvektoren erlaubt uns die Transformation von $\mathbf{x}_2$ auf Diagonalform
\[
\bm{\Lambda}_2=
\begin{pmatrix}7&0&0\\0&-2&0\\0&0&-2\end{pmatrix}.
\]

Eine ganz andere Situation liegt bei $\mathbf{M}_3$ vor. Wir berechnen die Eigenwerte und stellen fest, dass $\lambda_1=\lambda_2=\lambda_3=-1$.
Die Eigenvektoren berechnen wir aus  $(\mathbf{M}_3+\mathds{1})\cdot\mathbf{x}_i=0$. Wir stellen fest, dass $\text{Rg}(\mathbf{M}_3+\mathds{1})=2$ nicht einen Rangabfall von 3 aufweist, wie es bei einem dreifachen Eigenwert der Fall sein müsste. Also haben wir nur eine zweiparametrische Lösung vorliegen, die wir aus einer der linear abhängigen Gleichungen $(\mathbf{M}_3+\mathds{1})\cdot\mathbf{x}_i=0$ berechnen können:
\[
\begin{pmatrix}7&-6&5\\ 14&-12&10\\ 7&-6&5\end{pmatrix}
\begin{pmatrix}x_{i1}\\ x_{i2}\\ x_{i3}\end{pmatrix}=
\begin{pmatrix}0\\ 0\\ 0\end{pmatrix}
\]
Wir sehen sofort, dass diese drei Gleichungen linear abhängig sind 
\[ 7x_{i1}-6x_{i2}+5x_{i3}=0\]
und wir nur eine zweiparametrische Lösung bekommen. Wir wählen schlauerweise $x_{i1}=5\cdot s_i$ und $x_{i2}=5\cdot t_i$ und erhalten damit $x_{i3}\cdot s_i+6\cdot t_i$. Für $s_1=\frac{1}{5}$ und $t_1=\frac{2}{5}$ ergibt sich der Eigenvektor $\mathbf{x}_1=(1,2,1)^T$. Mit $s_2=\frac{3}{5}$ und $t_2=\frac{1}{5}$ ergibt sich der Eigenvektor $\mathbf{x}_2=(3,1,-3)^T$. Alle weiteren Eigenvektoren sind Linearkombination aus $\mathbf{x}_1$ und $\mathbf{x}_2$. D.h. das LGS gibt keine weiteren Eigenvektoren her. Anders gesagt $\mathbf{P}_3$ wäre singulär und nicht invertierbar und damit können wir $\mathbf{M}_3$ nicht auf Diagonalform bringen. Berechnen wir anstelle eines linear abhängigen dritten Eigenvektors den Lösungsvektor von $(\mathbf{M}_3+\mathds{1})^2\cdot\mathbf{x}_i=0$, dann lässt sich zeigen, dass ein $\mathbf{P}_3$ gebildet aus den ersten beiden Eigenvektoren und disem Lösungsvektor $\mathbf{M}_3$ zumidest auf die JORDANsche Normalform transformieren kann. Dies ist keine Diagonalform. Auf der Diagonalen stehen zwar die Eigenwerte, aber die obere Nebendiagonale ist mit Einsen besetzt. Und zwar stehen bei k-fachem eigenwert und Rangabfall $r\le k$ genau $k-r$ Einsen auf der oberen Nebendiagonale. Diese Aussage beweisen wir nicht. Berechnen wir zunächst den Lösungsvektor von $(\mathbf{M}_3+\mathds{1})^2\cdot\mathbf{x}_i=0$. Wir sehen, dass $(\mathbf{M}_3+\mathds{1})^2$ die Nullmatrix ist und damit die Lösung des Gleichungssystems beliebig. Ganz beliebig geht natürlich nicht, der Vektor muss zumindest linear unabhängig von den ersten beiden Eigenvektoren sein. Wir wählen $\mathbf{x}_3=(0,1,1)^T$ und konstruiren uns die Transformationsmatrix
\[
\mathbf{P}_3=
\begin{pmatrix}-1&0&3\\ -2&1&1\\ -1&1&-3\end{pmatrix}\quad\text{ wobei }\quad
\mathbf{P}_3^{-1}=
\begin{pmatrix}-4&3&-3\\ -7&6&-5\\ -1&1&-1\end{pmatrix}.
\]
Damit erhalten wir 
\[
\mathbf{P}_3^{-1}\mathbf{M}_3\mathbf{P}_3=
\begin{pmatrix}-1&1&0\\ 0&-1&0\\ 0&0&-1\end{pmatrix}
\]
\section{Das Computer Algebra System Maxima}
Maxima ist ein Computer Algebra System (CAS). Der Quellcode wurde seit 1982 von William Shelter\footnote{William Frederick Schelter (1947 – July 30, 2001) war Professor für Mathematik der  University of Texas in Austin. Er war Lisp Entwickler und  Programmierer (siehe \url{https://en.wikipedia.org/wiki/Bill_Schelter}).}  entwickelt und 1998 unter der GNU General Public License (GPL) veröffentlicht.

Maxima ist die Engine, die hinter den Stackaufgaben auf ILIAS steht, die wir im Rahmen dieses Kurses bearbeiten. Es ist daher ratsam, sich zumindest mit der Syntax von Maxima zu beschäftigen, da diese zur Eingabe bei den Aufgaben benutzt werden muss. Die Syntax ist einigermaßen intuitiv, es schadet aber nicht, sich einmal damit zu beschäftigen. Weitere Information ist hier 
\url{http://maxima.sourceforge.net/docs/manual/de/maxima.html#SEC_Top} 
zu finden.
\newpage
\chapter{Gewöhnliche Differentialgleichungen}
Die allgemeine Form einer Differentialgleichung sei folgendermaßen gegeben
\begin{equation}
  \frac{d\mbf{y}(t)}{dt}=\mbf{f}(t,\mbf{y}(t))
  \label{eq:DGLallgemein}
\end{equation}
Anstatt den Differentialoperator $\frac{d}{dt}$ auszuschreiben, benutzen wir
für (\ref{eq:DGLallgemein}) auch die Schreibweise
$\dot{\mbf{y}}=f(t,\mbf{y}(t))$.  Dabei bezeichne $t$ die unabhängige
Veränderliche, $\mbf{f}(t,\mbf{y})$ einen Vektor bekannter Funktion und
$\mbf{y}$ einen Vektor gesuchter Funktionen $y_i(t)$.  Wir betrachten
verschiedene Differentialgleichungen, benennen deren Eigenschaften und erkennen
die verschiedenen Typen:
\begin{enumerate}
	\item Linear und nichtlinear, wie z.B.\
	  \[m\ddot{x}(t)+c\dot{x}(t)+kx(t)=f(t),\]
	  eine lineare inhomogene Differentialgleichung 2. Ordnung,
	  die den gedämpften und getriebenen harmonischen Oszillator bechreibt, während
	\[\frac{d^2x(t)}{dt^2}+\mu(x(t)^2-1)\frac{dx(t)}{dt}+x(t)= 0 \]
	  eine nichtlineare Bewegungsgleichung für $x(t)$ ist. Sie beschreibt den so genannten
	  van der Pol Oszillator.  
	\item erste, zweite und höhere Ordnung 
	  \[ \dot{x}(t)+\frac{1}{\tau}x=h(t)\] 
	  beschreibt einen getriebenen Relaxator, den wir gleich näher
	  betrachten werden. Die rechte Seite dieser Gleichung bezeichnen wir
	  als Inhomogenität.
		%- in der Systemtheorievbezeichneten wir damit ein
		%Verzögerungsglied 1.\ Ordnung ($\mbox{PT}_1$-Glied). 
	  Ein Beispiel für eine Gleichung 2.\ Ordnung ist entweder der
	  van der Pol Oszillator oder der Harmonische Oszillator (siehe
	  oben). 
		% Finde ein entsprechendes elementares zeitinvariantes
		% \"Ubertragungsglied (Hinweis: Skript Systemtheorie)!
	\item System von DEs  1.\ Ordnung 
		\begin{eqnarray*}
			\frac{dx}{dt} &=& x(\alpha - \beta y) \\
			\frac{dy}{dt} &=& - y(\gamma - \delta x) 
		\end{eqnarray*}
		die bekannte Räuber-Beute-Gleichungen oder auch Lotka-Volterra-Gleichungen.
	\item Transformation von Gleichungen höherer Ordnung auf ein System von Gleichungen 1.\ Ordnung.
		Gegeben der gedämpfte Harmonische Oszillator
		\[m\ddot{x}(t)+c\dot{x}(t)+kx=f(t)\]. Substituiere
		$y_1=x$ und $y_2=\dot{x}$ und wir erhalten zwei Gleichungen erster
		Ordnung anstatt der ursprünglichen Gleichung zweiter Ordnung, nämlich 
		\begin{eqnarray*} 
		  \dot{y_1} &=& y_2\\
	 	  m\dot{y_2}&=&-k\cdot y_1-c\cdot y_2+f(t)
	       \end{eqnarray*}
	     \item Der Grad einer Differentialgleichung ist der Exponent der Potenz der höchsten vorkommenden Ableitung.
\end{enumerate}
Bei all diesen Differentialgleichungen sind wir immer an einer Lösung für
einen bestimmten Anfangswert interessiert, also z.B.\ $x(t=0)=x_0$ etc. Um ein
System von Gleichungen lösen zu können müssen wir für jede Variable
Anfangsbedingungen angeben. Da aus einer Gleichung n-ter Ordnung ein System mit
n Variablen wird, schließen wir daraus, dass wir für eine solche
Differentialgleichung ebenfalls n Anfangsbedingungen brauchen.
%
\section{Allgemeine analytische Lösung}\label{sec:analyticsolu}
Wenn wir gewöhnliche Differentialgleichungen analytisch lösen wollen so ist
dies nicht in allen Fällen möglich. Betrachten wir das Anfangswertproblem
für die gewöhnliche Differentialgleichung
\[ \dot{y}(t) = f\left(t,y(t)\right)\]
mit $y(t_0)=y_0$. Die Funktion $f(t,y(t))$ sei stetig im Bereich der
$(t,y)$-Ebene, gegeben durch $t_0-a\le t\le t_0+a$ und $y_0 -b\le y\le y_0+b$.
Der Nachweis der Existenz und Eindeutigkeit einer Lösung ist in diesem Falle
nicht immer möglich. Wir wollen eine alternative Vorgehensweise vorschlagen.

Wir nehmen an $y(t)$ sei eine Lösung des Anfangswertproblems und es gelte
$y_0 -b\le y\le y_0+b$ für $t\in [t_0-\alpha,t_0+\alpha]$ mit $0<\alpha\le a$,
so ergibt sich aus obiger Voraussetzung für $f(t,y(t))$, dass die Lösung im
Intervall $[t_0-\alpha,t_0+\alpha]$ eine stetige Ableitungsfunktion
$\dot{y}(t)$ besitzt.  Also lässt sich die Lösungsfunktion dort als
Stammfunktion dieser Ableitung darstellen 
\begin{equation}\label{eq:Integral}
  y(t)=y_0+\int_{t_0}^t f(s,y(s))ds 
\end{equation}
Wobei $t_0-a\le t\le t_0+a$ gilt. Dies bedeutet aber, dass jede Lösung des
obigen Anfangswertproblems auch Lösung der Integralgleichung
(\ref{eq:Integral}) ist. Umgekehrt ist auch jede Lösung von
(\ref{eq:Integral}) Lösung des Anfangswertproblems.
\section{Lösung durch Iteration}
Schreiben wir die Integralgleichung (\ref{eq:Integral}) in der Form
\begin{equation}
  y(t)=\Phi[y]\mbox{ mit } \Phi[y]=y_0+\int_{t_0}^t f(s,y(s))ds
  \label{eq:Iteration}
\end{equation}
\begin{note}{Funktional statt Integral}
  In (\ref{eq:Iteration}) schreiben wir absichtlich $\Phi[y]$, d.h. das $y$ in
  eckigen Klammern geschrieben, weil $\Phi[y]$ ein Funktional ist. Es hängt
  nicht nur vom Zeitpunkt $t$ in der unabhängigen Variablen ab, sondern vom
  gesamten Verlauf der Funktion $y(t)$. Wir werden im Kapitel \ref{chap:Variationsrechnung}
  Variationsrechnung darauf zurückkommen. 
\end{note}
Wir sehen, dass die Lösung von (\ref{eq:Iteration}) wie der Fixpunkt $y^*$ der
Gleichung \[ y=\varphi(y)\] zu bestimmen ist. Hierbei gehen wir iterativ vor.
Wir starten bei einem Rohwert $y_0$ und berechnen sukzessiv nach der Vorschrift
\[ y_{n+1} = \varphi(y_n) \] 
Jetzt müssen wir nur noch zeigen, dass die Iteration konvergiert. Dies ist für
die Lösung der Integralgleichung (\ref{eq:Integral}) der Fall, wenn die
Abbildung $\phi$ kontrahierend ist.
\begin{example}{Iterative Lösung}
Gegeben die Differentialgleichung
\[
  \frac{d}{dt}y(t)=-\alpha y(t) 
\]
Zeige, dass die iterative Lösung im Limes mit der analytischen Lösung
übereinstimmt.
\end{example}
\section{Gleichungen erster Ordnung und ersten Grades}
Wir gehen von einer Differentialgleichung erster Ordnung und ersten Grades 
\[   p(t,y)+q(t,y)\dot{y}(t)=0 \]
aus, die sich in der Form
\begin{equation}\label{eq:DGLVarKoeff}
  p(t,y)dt+q(t,y)dy=0
\end{equation}
schreiben lässt. Wobei $p(t,y)$ und $q(t,y)$ den gemeinsamen Definitionsbereich
$a\le t\le b$ und $\alpha\le y\le\beta$ haben
\begin{example}{Umformung}
  \[\frac{dy}{dt}+\frac{y+t}{y-t}=0\]
  kann mit $p(t,y)=y+t$ und $q(t,y)=y-t$ auf die Form
  \[(y+t)dt+(y-t)dy=0\]
  gebracht werden.
\end{example}
\subsection{Separierbare Differentialgleichungen}
Den Spezialfall
\begin{equation}
  p(t)+q(y)\dot{y}(t)=0
  \label{eq:getrennteVar}
\end{equation}
nennen wir eine {\it Differentialgleichung mit getrennten Variablen}.

Existieren für $p$ und $q$ im angegebenen Definitionsbereich die Integrale $P$
und $Q$, mit $\dot{P}(t)=p(t)$ und $\frac{dQ(y)}{dy}=q(y)$ so ist
\[\mu(t,y)=P(t)+Q(y)=const.\]
denn es ist wegen (\ref{eq:getrennteVar})
\[ d\mu=\dot{P}(t)dt+\frac{dQ(y)}{dy}dy=p(t)dt+q(y)dy=0\]
\begin{example}{Separierbare Differentialgleichung}
a) $t\dot{y}=y$ und b) $y\dot{y}=-t$
\begin{description}
  \item[a)] Für $t\ne 0$ und $y\ne =0$ erhalten wir 
  \[\frac{dy}{y}=\frac{dt}{t}\]
\item[b)] $ydy=-tdt$
\end{description}
\end{example}
\subsection{Exakte Differentialgleichungen}
Ist $p(t,y)dt+q(t,y)dy$ das vollständige Differential einer Funktion
$\mu(y,t)$, d.h.\ ist 
\[d\mu(y,t)=\frac{\partial\mu(y,t)}{\partial t}dt+
\frac{\partial\mu(y,t)}{\partial y}dy=p(t,y)dt+q(t,y)dy\]
dann nennen wir (\ref{eq:DGLVarKoeff}) eine {\it exakte Differentialgleichung}
und $\mu(y,t)=const.$ ist eine allgemeine Lösung. Damit (\ref{eq:DGLVarKoeff})
eine exakte Differentialgleichung ist, muss die Bedingung
\begin{equation}
  \frac{\partial p(t,y)}{\partial y}=  \frac{\partial q(t,y)}{\partial t}
  \label{eq:BedExaktheit}
\end{equation}
erfüllt sein.

Nicht immer sieht man allerdings der Differentialgleichung die Exaktheit an.
Selbst wenn wir sehen, dass (\ref{eq:DGLVarKoeff}) keine exakte
Differentialgleichung ist, da (\ref{eq:BedExaktheit}) nicht erfüllt ist, gibt
es aber unter Umständen die Möglichkeit eine Funktion zu finden, mit der man
die Gleichung multiplizieren kann, sodass diese exakt wird. 
\begin{note}{Integrierender Faktor}
  $3ydt+2tdy=0$ ist keine exakte Differentialgleichung. Multipliziert man aber
  $\Xi(y,t)=t^2y$ mit dieser Gleichung, so erhält man $3t^2y^2dt+2t^3ydy=0$.
  Dies kann man als vollständiges Differential der unktion $\mu(y,t)=t^3y^2$
  schreiben und somit ist $t^3y^2=C$ einen allgemeine Lösung der
  Differentialgleichung. Die Funktion $\Xi(y,t)$ heisst integrierender Faktor.
\end{note}
Angenommen (\ref{eq:DGLVarKoeff}) sei keine exakte Differentialgleichung, 
dann müssen wir einen solchen versuchen zu finden
\begin{enumerate}
    \item Wenn 
    \[\frac{\frac{\partial p(t,y)}{\partial y}-\frac{\partial q(t,y)}{\partial t}}{q(t,y)}=f(t)\] 
    eine Funktion von t allein ist, dann ist $e^{\int f(t)dt}$ ein integrierender Faktor. 
    Ebenso wenn
    \[\frac{\frac{\partial p(t,y)}{\partial y}-\frac{\partial q(t,y)}{\partial t}}{p(t,y)}=-g(y)\]
    nur eine Funktion von y ist, dann ist $e^{\int g(y)dy}$ ein integrierender Faktor.
    \item Ist (\ref{eq:DGLVarKoeff}) homogen, d.h. sind $p(t,y)$ und $q(t,y)$
      homogene Funktionen vom selben Grad, und $p(t,y)\cdot t+q(t,y)\cdot y\ne
      0$, dann ist $\frac{1}{p(t,y)\cdot t+q(t,y)\cdot y}$ ein integrierender
      Faktor.
    \item Wenn (\ref{eq:DGLVarKoeff}) in der Form $y\cdot f(t\cdot y)dt+t\cdot g(t\cdot y)dy=0$
    geschrieben werden kann, dann ist $\frac{1}{t\cdot y(f(t\cdot y)-g(t\cdot y))}=\frac{1}{p\cdot t-q\cdot y}$ ein integrierender Faktor.
    \item Durch genaues Hinschauen kann man durch Umgruppieren von Termen einen
      integrierenden Faktor finden, wenn man bestimmte Gruppen von Termen als
      Teil eines vollständigen Differentials identifiziert.
\end{enumerate}
\Comment{
%\begin{example}{}
Beispiele zu Gruppe von Termen, Integrierender Faktor, Vollständiges Differential.
%\end{example}
} 
\chapter{Gewöhnliche lineare homogene Differentialgleichungen}
\Comment{2. Vorlesung}
\section{Gleichungen mit konstanten Koeffizienten}
Eine lineare Differentialgleichung n-ter Ordnung ist ein Zusammenhang folgender
Art
\begin{equation}\label{eq:DGLOrdnungN}
a_n\frac{d^n y(t)}{dt^n}+a_{n-1}\frac{d^{n-1} y(t)}{dt^{n-1}}+
\cdots +a_{1}\frac{d y(t)}{dt}+a_0y(t)=f(t)
\end{equation}
Dabei ist $y(t)$ eine Funktion der unabhängigen Veränderlichen $t$.
Differentialgleichungen vom Typ (\ref{eq:DGLOrdnungN}) beschreiben zahlreiche
Phänomene in Naturwissenschaft und Technik. Zunächst einmal wollen wir die
linearen Differentialgleichungen mit konstanten Koeffizienten untersuchen. Die Funktion $f(t)$ sei nur von der unabhängigen Veränderlichen $t$ ab\-hängig. Wenn $f(t)=0$ nennen wir (\ref{eq:DGLOrdnungN}) eine homogene Differentialgleichung, wenn dem nicht so ist, heisst (\ref{eq:DGLOrdnungN}) inhomogen.
\begin{note}{Homogene Funktion vom Grad n}
  Eine Funktion $f(x)$ wird homogen vom Grad $n$ genannt, wenn gilt $f(\lambda\cdot x)=\lambda^n\cdot f(x)$. Deshalb nennt man $f(t)$ der linearen Differentialgleichung (\ref{eq:DGLOrdnungN}) auch die \textit{Inhomogenität} der Differentialgleichung.
\end{note}

Man beachte, dass die Koeffizienten der Ableitungen von $y(t)$ nicht
notwendigerweise Konstanten sein müssen, sondern bekannte Funktionen $a_k(t)$
sein können. Wichtig für die Bezeichnung als lineare Differentialgleichung ist
ganz einfach die Tatsache, dass das Superpositionsprinzip gilt. 
\begin{note}{Lineare Superposition von Lösungen}
D.h. wenn $y_a(t)$ und $y_b(t)$ je eine Lösung der homogenen linearen
Differentialgleichung sind, dann ist auch $y(t)=\alpha y_a(t) + \beta y_b(t)$
eine Lösung. Es löse $y_f(t)$ die Differentialgleichung mit der Inhomogenität
$f(t)$ und $y_g(t)$ diejenige mit der Inhomogenität $g(t)$, dann löst
$y(t)=\alpha y_f(t) + \beta y_g(t)$ die Differentialgleichung mit der
Inhomogenität $h(t)=\alpha f(t) + \beta g(t)$. Man beweise diese Aussage!
\end{note}
%
\section{Lösung der Gleichungen mit konstanten Koeffizienten}
Zur Lösung der Gleichungen vom Typ (\ref{eq:DGLOrdnungN}) haben wir
verschiedene Herangeshensweisen. Dabei gehen wir davon aus, dass es genau eine
Lösung von (\ref{eq:DGLOrdnungN}) mit $f(t)=0$, also der homogenen
Differentialgleichung, gibt, die den Bedingungen 
\begin{equation}
y(t_0)=b_0,\,\dot{y}(t_0)=b_1,\,\dots y^{(n-1)}(t_0)=b_{n-1}
  \label{eq:ICsGDGL}
\end{equation}
genügt. Man möge den Beweis dieses Satzes in den einschlägigen Lehrbüchern
nachlesen.
\section{Taylorreihe und  Potentzreihenansatz}
Zur Lösung von (\ref{eq:DGLOrdnungN}) liegt es nahe, den Ansatz
\begin{equation}
  y(t)=\sum\limits_{\nu=0}^{\infty}c_\nu t^\nu
  \label{eq:AnsatzPotenz}
\end{equation}
zu versuchen. D.h. wir vermuten die Lösung in der Form (\ref{eq:AnsatzPotenz})
finden zu können. Hierzu gibt es zwei Herangehensweisen. Einmal mit Hilfe der
Taylorreihe und zum anderen durch direktes Einsetzen von
(\ref{eq:AnsatzPotenz}). 

Wir arbeiten beides am Beispiel des Relaxators aus.
\begin{example}{Taylorreihe verglichen mit Potenzreihenansatz}
  Gegeben die homogene Differentialgleichung $\dot{y}(t)+a_0y(t)=0$ mit
  Anfangsbedingung $y(0)=b_0$ - wir setzen o.B.d.A. $t_0=0$. Wir haben
  \[ \dot{y}(t)=-a_0y(t),\, \ddot{y}(t)=-a_0\dot{y}(t),\,\dots\, 
      y^{(n)}(t)=-a_0y^{(n-1)}(t)\]
  Damit erhalten wir die Ableitungen bei $t=0$
  \[ y(0)=b_0,\, \dot{y}(0)=-a_0b_0,\, \ddot{y}(0)=a_0^2b_0,\,\dots\, 
     y^{(n)}(0)=(-1)^na_0^nb_0\]
  und können die Taylorreihe für $y(t)$ um $t=0$ schreiben als
  \begin{equation}
    y(t)=b_0\sum\limits_{\nu=0}^{\infty}(-1)^\nu\frac{a_0^\nu}{\nu!}t^\nu=b_0e^{-a_0t}
    \label{eq:SoluTaylor}
  \end{equation}
Wenn wir (\ref{eq:AnsatzPotenz}) direkt einsetzen, dann erhalten wir die Lösung
durch Koeffizientenvergleich.
  \[ \sum\limits_{\nu=0}^{\infty}\nu c_\nu t^{\nu-1}+ 
      a_0\sum\limits_{\nu=0}^{\infty}c_\nu t^\nu=
      \sum\limits_{\mu=0}^{\infty}(\mu+1)c_{\mu+1}t^\mu+ 
      a_0\sum\limits_{\nu=0}^{\infty}c_\nu t^\nu=
      \sum\limits_{\nu=0}^{\infty} \left((\nu+1)c_{\nu+1}+a_0c_\nu\right) t^\nu
  \]
  Aus der Anfangsbedingung erhalten wir eine Rekursionsformel für die $c_\nu$
  \begin{equation}
    (\nu+1)c_{\nu+1}+a_0c_\nu=0
    \label{eq:Rekursion}
  \end{equation}
  Mit $y(0)=b0$ haben wir $c_0=b_0$ und finden damit die $c_\nu$ aus
  (\ref{eq:Rekursion}) 
    \begin{equation}
      c_{\nu}=(-1)^\nu\frac{a_0^\nu b_0}{\nu!}\qquad (\nu=0,1,2,\dots).
    \label{eq:PotenzKoeffizienten}
  \end{equation}
  Mit den Koeffizienten (\ref{eq:PotenzKoeffizienten}) erhalten wir die Lösung
  \begin{equation}
    y(t)=\sum\limits_{\nu=0}^{\infty}(-1)^\nu\frac{a_0^\nu b_0}{\nu!}t^\nu=b_0e^{-a_0t}
    \label{eq:SoluPotenz}
  \end{equation}
\end{example}
%
\section{Der Lösungsansatz $y(t)=e^{\lambda t}$}
Setzen wir $y(t)=e^{\lambda t}$ in (\ref{eq:DGLOrdnungN}) ein, so erhalten wir
\begin{equation}
  \lambda^n+a_{n-1}\lambda^{n-1}+\dots +a_0=0
  \label{eq:Charakteristische}
\end{equation}
Die Gleichung (\ref{eq:Charakteristische}) hat $n$ Nullstellen und somit haben
wir auch $n$ Lösungen, die wir als Superposition zur Lösung des
Anfangswertproblems angeben können. Für unseren Relaxator bedeutet dies
\[ \lambda e^{\lambda t}+a_0e^{\lambda t}=0\]
Und daraus erhalten wir die Lösung $y(t)=ce^{-a_0 t}$. Um die Anfangsbedingung
zu erfüllen, setzen wir $c=b_0$.

\begin{example}{Der Exponentialansatz für eine Gleichung 2. Ordnung}
  Gegeben die Gleichung $\ddot{y}(t)+2\dot{y}(t)-3y(t)=0$ mit den
  Anfangsbedingungen $y(0)=2$ und $\dot{y}(0)=-2$. Die charakteristische
  Gleichung (\ref{eq:Charakteristische}) lautet für diesen Fall
  $\lambda^2+2\lambda-3=0$ mit den beiden Lösungen $\lambda_1=1$ und
  $\lambda_2=-3$. Wir schreiben die Lösung der Gleichung als Superposition
  $y(t)=c_1e^{\lambda_1 t}+c_2e^{\lambda_2 t}=c_1e^{t}+c_2e^{-3 t}$. Wenn wir
  die beiden Anfangsbedingunen erfüllen, dann sind $c_1=c_2=1$ festgelegt.
  Damit ist die Lösung $y(t)=e^t+e^{-3t}$.
\end{example}
Damit erhebt sich die berechtigte Frage, warum man den Potenzreihenansatz dann
überhaupt in Betracht zieht. Wir erörten dies an folgendem Beispiel:
\begin{example}{Mehrfachlösungen des charakteristishen Polynoms}
  \[\ddot{y}(t)-2\dot{y}(t)+y(t)=0,\mbox{ mit } y(0)=1,\quad \dot{y}(0)=2\]
Mit dem Exponentialansatz $e^{\lambda t}$ erhalten wir das charakteristische
Polynom $\lambda^2-2\lambda+1=(\lambda -1)^2=0$. Dieses hat eine
doppelte Nullstelle, aber wie lautet die zweite Lösung?
\end{example}
Versuchen wir es mit dem Potenzreihenansatz (\ref{eq:AnsatzPotenz}) wir
erhalten die Rekursionsformel
\[ (\nu+1)(\nu+2)c_{\nu+2}-2(\nu+1)c_{\nu+1}+c_\nu=0\]
Mit $d_\nu=c_{\nu}\nu!$ erhalten wir die Rekursionsformel
$d_{\nu+2}-2d_{\nu+1}+d_\nu=0$, wobei wegen der Anfangsbedingungen $d_0=c_0=1$
und $d_1=c_1=2$ sein muss und wir sehen dass $d_\nu=\nu+1$. Damit ist aber
\[ 
  y(t)=\sum\limits_{\nu=0}^{\infty}\frac{1+\nu}{\nu!}t^\nu=
        \sum\limits_{\nu=0}^{\infty}\frac{t^\nu}{\nu!}
       +\sum\limits_{\nu=1}^{\infty}\frac{t^\nu}{(\nu-1)!}=
        e^t+t\sum\limits_{\mu=0}^{\infty}\frac{t^\mu}{\mu!}=e^t+te^t
\]       
Dies legt die Vermutung nahe, dass bei einer r-fachen Lösung $\lambda_1$ des
charakteristischen Polynoms neben $e^{\lambda_1 t}$ auch $te^{\lambda_1
t},\dots,t^{r-1}e^{\lambda_1 t}$ Lösungen sind. Man zeige dies durch einsetzen
der einzelnen Lösungen in die homogene Differentialgleichung
$(D-\lambda_1)^r[y(t)]=0$!
%
\section{Das Fundamentalsystem von Lösungen}
Die Frage nach der Zahl der Lösungen soll im folgenden behandelt werden, d.h.
wir suchen das Fundamentalsystem von Lösungen zu (\ref{eq:DGLOrdnungN}).
\begin{example}{Fundamentalsystem der Gleichung 2. Ordnung}
  Wir wissen, dass für die Gleichung 2. Ordnung 
  \[\ddot{y}(t)+a_0y(t)=0 \]
  zwei Bedingungen vorliegen müssen, damit die Lösung eindeutig angegebn werden
  kann. Also nehmen wir die Bedingunen
  \begin{align}
  y(0)=&b_0\mbox{ und }\dot{y}(0)=0\mbox{ oder}\nonumber\\
  y(0)=&0\mbox{ und }\dot{y}(0)=b_1\nonumber\\
  \end{align}
  Der Potenzreihenansatz
  \[ y(t)=\sum\limits_{\nu=0}^\infty c_\nu t^\nu \]
  führt auf
  \begin{align*}
  \sum\limits_{\nu=2}^\infty\nu(\nu-1)c_\nu t^{\nu-2}+a_0\sum\limits_{\nu=0}^\infty c_\nu t^\nu &=0\\
  \sum\limits_{\nu=0}^\infty\left( (\nu+2)(\nu+1)c_{\nu+2} +a_0c_\nu\right) t^\nu &=0
  \end{align*}
  und mit der ersten Anfangsbedingung auf die Lösungen $y_1(t)$ mit
  \[ c_{2\mu}=(-1)^\mu\frac{a_0^\mu b_0}{(2\mu)!},\quad c_{2\mu+1}=0\]
  während die zweite Anfangsbedingung auf die Lösung $y_2$ mit
  \[ c_{2\mu}=0, \quad c_{2\mu+1}=(-1)^\mu\frac{a_0^\mu b_1}{(2\mu+1)!}\]
  führt. Damit erfüllt die Summe der beiden Lösungen die allgemeine
  Anfangsbedingung
  \[ y(0)=b_0\mbox{ und }\dot{y}(0)=b_1\]
\end{example}
Das obige Beispiel legt nahe, bei homogenen Differentiagleichungen n-ter
Ordnung, wie in (\ref{eq:DGLOrdnungN}) mit $f(t)=0$ angegeben, ähnlich
vorzugehen, indem wir der Reihe nach folgende Anfangsbedingungen aufstellen
\begin{align*}
  y(t_0)=1,\dot{y}(t_0)=0,&\dots\dots\dots,y^{(n-1)}(t_0)=0;\\
  y(t_0)=0,\dot{y}(t_0)=1,&\dots\dots\dots,y^{(n-1)}(t_0)=0;\\
  &\dots\dots\dots \\
  y(t_0)=0,\dot{y}(t_0)=0,&\dots\dots\dots,y^{(n-1)}(t_0)=1;
  \label{eq:NICs}
\end{align*}
Wir bezeichnen die entsprechenden Lösungen mit $y_1(t),y_2(t),\dots,y_n(t)$.
Für die Anfangsbedingung
$y(t_0)=b_0,\dot{y}(t_0)=b_1,\dots,y^{(n-1)}(t_0)=b_{n-1}$ können wir damit
sofort die Lösung in der Form
\begin{equation}
  y(t)=b_0y_1(t)+b_1y_2(t)+\dots+b_{n-1}y_n(t)
  \label{eq:FundamentalSuperpose}
\end{equation}
angeben. Dies ist allerdings bisher auf die Stelle $t_0$ beschränkt, wir wollen
versuchen eine allgemeine Aussage zu machen. Wir bilden mit den obigen Lösungen 
\[y(t)=c_1y_1(t)+c_2y_2(t)+\dots+c_ny_n(t)\qquad (c_1,c_2,\dots,c_n\in\mathbb{R})\]
was wieder Lösung der Differentialgleichung (\ref{eq:DGLOrdnungN}) ist.

Das lineare Gleichungssystem
\begin{align}
c_1y_1(t_0)+c_2y_2(t_0)&+{\dots\dots\dots}+c_ny_n(t_0)=b_0\nonumber\\
c_1\dot{y}_1(t_0)+c_2\dot{y}_2(t_0)&+{\dots\dots\dots}+c_n\dot{y}_n(t_0)=b_1\nonumber\\
c_1\ddot{y}_1(t_0)+c_2\ddot{y}_2(t_0)&+{\dots\dots\dots}+c_n\ddot{y}_n(t_0)=b_2\nonumber\\
  &\dots\dots\dots \\
  c_1y^{(n-2)}_1(t_0)+c_2y^{(n-2)}_2(t_0)&+{\dots\dots\dots}+c_ny^{(n-2)}_n(t_0)=b_{n-2}\nonumber\\
  c_1y^{(n-1)}_1(t_0)+c_2y^{(n-1)}_2(t_0)&+{\dots\dots\dots}+c_ny^{(n-1)}_n(t_0)=b_{n-1}\nonumber
  \label{eq:LGSatt0}
\end{align}
hat die eindeutige Lösung $c_1=b_0$, $c_2=b_1$,\dots, $c_n=b_{n-1}$. Das
bedeutet dass die Koeffizientendeterminante
\begin{equation}
  W(t_0)=\begin{vmatrix}
           y_1(t_0)&y_2(t_0)&\quad\dots\quad &y_n(t_0)\\
           \dot{y}_1(t_0)&\dot{y}_2(t_0)&\quad\dots\quad &\dot{y}_n(t_0)\\
           \ddot{y}_1(t_0)&\ddot{y}_2(t_0)&\quad\ddots\quad &\ddot{y}_n(t_0)\\
	   \dots&\dots &\quad\dots\quad&\dots\\
           y^{(n-2)}_1(t_0)&y^{(n-2)}_2(t_0)&\quad\dots\quad &y^{(n-2)}_n(t_0)\\
           y^{(n-1)}_1(t_0)&y^{(n-1)}_2(t_0)&\quad\dots\quad &y^{(n-1)}_n(t_0)
         \end{vmatrix}
  \label{eq:Wronskiant0}
\end{equation}
von null verschieden ist. Die Determinante (\ref{eq:Wronskiant0}), wird
Wronskideterminante genannt und sie ist also, wenn eine Lösung für die
Anfangsbedingungen bei $t_0$ vorliegt, von null verschieden. Wenn dies aber für
beliebige Stellen $t$ der Fall sein soll, müssen wir das allgemeine Verhalten
der Wronskideterminante untersuchen, d.h. sie muss für jedes beliebige $t$ von
null verschieden sein.

Dazu bilden wir die Ableitung von $W(t)$ und erhalten
\begin{align}
  \dot{W}(t)=&\begin{vmatrix}
    \dot{y}_1(t)&\quad\dots\quad &\dot{y}_n(t)\\
           \dot{y}_1(t)&\quad\dots\quad &\dot{y}_n(t)\\
           \ddot{y}_1(t)&\quad\ddots\quad &\ddot{y}_n(t)\\
                       &\quad\dots\quad&\dots\\
           y^{(n-1)}_1(t)&\quad\dots\quad &y^{(n-1)}_n(t)
         \end{vmatrix}+
	 \begin{vmatrix}
           y_1(t)&\quad\dots\quad &y_n(t)\\
           \ddot{y}_1(t)&\quad\dots\quad &\ddot{y}_n(t)\\
           \ddot{y}_1(t)&\quad\dots\quad &\ddot{y}_n(t)\\
                       &\quad\dots\quad&\dots\\
           y^{(n-1)}_1(t)&\quad\dots\quad &y^{(n-1)}_n(t)
         \end{vmatrix}+\dots\\
      \dots+&\begin{vmatrix}
           y_1(t)&\quad\dots\quad &y_n(t)\\
           \dot{y}_1(t)&\quad\dots\quad &\dot{y}_n(t)\\
                       &\quad\dots\quad&\dots\\
           y^{(n-1)}_1(t)&\quad\dots\quad &y^{(n-1)}_n(t)\\
           y^{(n-1)}_1(t)&\quad\dots\quad &y^{(n-1)}_n(t)
         \end{vmatrix}+
      \begin{vmatrix}
           y_1(t)&\quad\dots\quad &y_n(t)\\
           \dot{y}_1(t)&\quad\dots\quad &\dot{y}_n(t)\\
                       &\quad\dots\quad&\dots\\
           y^{(n-2)}_1(t)&\quad\dots\quad &y^{(n-2)}_n(t)\\
           y^{(n)}_1(t)&\quad\dots\quad &y^{(n)}_n(t)
         \end{vmatrix}
  \label{eq:WronskianDerivative1}
\end{align}
Nur die letzte der resultierenden Determinanten ist von null verschieden.
O.b.d.A. nehmen wir $a_n=1$ in (\ref{eq:DGLOrdnungN}) an\footnote[1]{Wir teilen
(\ref{eq:DGLOrdnungN}) durch $a_n$, vorausgesetzt $a_n\ne0$, denn wenn das
nicht der Fall ist, liegt eine Differentialgleichung (n-1)-ter Ordnung vor.}
und erhalten damit
\begin{equation}
  \begin{vmatrix}
           y_1(t)&\quad\dots\quad &y_n(t)\\
           \dot{y}_1(t)&\quad\dots\quad &\dot{y}_n(t)\\
                       &\quad\dots\quad&\dots\\
      -a_{n-1}y^{(n-1)}_1(t)-\dots-a_0y_1&\quad\dots\quad 
      &-a_{n-1}y^{(n-1)}_n(t)-\dots-a_0y_n(t)
         \end{vmatrix}
  \label{eq:WronskianDerivative2}
\end{equation}
Nun multiplizieren wir die erste Zeile von (\ref{eq:WronskianDerivative2}) mit
$a_0$ und addieren sie zur letzten Zeile, sukzessive bis $a_{n-2}$
multipliziert mit den vorletzten Zeile und wir erhalten
\begin{equation}
  \dot{W}(t)=-a_{n-1}W(t)
  \label{eq:WronskianDGL}
\end{equation}
Für diese Differentialgleichung haben wir eine Anfangsbedinung $W(t_0)=w_0$ und
können sie damit lösen. Wir erhalten
\begin{equation}
  W(t)=w_0e^{-a_{n-1}(t-t_0)}
  \label{eq:WronskianGeneral}
\end{equation}
Da die Exponentialfunktion keine Nullstellen hat, schließen wir aus
(\ref{eq:WronskianGeneral}), dass wenn $W(t)$ an irgendeiner Stelle von null
verschieden ist, sie üeberall von null verschieden ist. Wir sehen also, da
$W(t)$ im gegebenen Fall nicht null werden kann, dass immer eine Lösung
existiert.
\begin{note}{}
  Beachte das über multiple Lösungen des charakteristischen Polynoms oben
  gesagte!
\end{note}
\chapter{Lösung der inhomogenen Gleichung}
\section{Variation der Konstanten}
Diese Methode lässt sich auch bei variablen Koeffizienten anwenden. Die Lösung
der homogenen Gleichung hierfür, soll allerdings erst später besprochen werden.
Daher gehen wir von einer Differentialgleichung der Form (\ref{eq:DGLOrdnungN})
aus. Die Idee hinter der Methode ist die gefundenen Fundamentallösungen zu
superponieren, die benutzten Koeffizienten aber als noch zu bestimmende
Funktionen anzusetzen. Wir setzen also folgendes an
\begin{equation}
  y(t)=c_1(t)y_1(t)+c_2(t)y_2(t)+\cdots+c_n(t)y_n(t)
  \label{eq:AnsatzVariationKonstanten}
\end{equation}
Nun brauchen wir natürlich auch $n$ Bedingungen um die Funktionen $c_i(t)$ zu
bestimmen. Wir nehmen an dass die $c_i(t)$ mindestens einmal differenzierbar seien.
Die Lösung $y(t)$ hingegen muss n-mal differenzierbar sein. Dann legen wir folgendes fest für die $c_i(t)$
\begin{align}
  \dot{y}(t)=&c_1(t)\dot{y}_1(t)+c_2(t)\dot{y}_2(t)+\cdots+c_n(t)\dot{y}_n(t)\nonumber\\
  &+\underbrace{\dot{c}_1(t)y_1(t)+\dot{c}_2(t)y_2(t)+\cdots+\dot{c}_n(t)y_n(t)}_{=0}\nonumber\\
  \ddot{y}(t)=&c_1(t)\ddot{y}_1(t)+c_2(t)\ddot{y}_2(t)+\cdots+c_n(t)\ddot{y}_n(t)\nonumber\\
  &+\underbrace{\dot{c}_1(t)\dot{y}_1(t)+\dot{c}_2(t)\dot{y}_2(t)+\cdots
               +\dot{c}_n(t)\dot{y}_n(t)}_{=0}\nonumber\\
	       &\cdots\cdots\cdots\nonumber\\
	       y^{(n)}(t)=&c_1(t)y^{(n)}_1(t)+c_2(t)y^{(n)}_2(t)+\cdots+c_n(t)y^{(n)}_n(t)\nonumber\\
  &+\underbrace{\dot{c}_1(t)y^{(n-1)}_1(t)+\dot{c}_2(t)y^{(n-1)}_2(t)+\cdots
  +\dot{c}_n(t)y^{(n-1)}_n(t)}_{=g(t)}\label{eq:Diffbarkeitci}
\end{align}
Gegeben, dass diese Bedingungen erfüllt sind, so ist
(\ref{eq:AnsatzVariationKonstanten}) Lösung der Differentialgleichung. Dies
können wir sofort erkennen, denn wenn wir die so berechneten Ableitungen aus
(\ref{eq:Diffbarkeitci}) in die Differentialgleichung einsetzen, ist diese
erfüllt.

Für die $\dot{c}_k(t)$ dagegen habe wir das Gleichungssystem
\begin{align}
  \dot{c}_1(t)y_1(t)&+\dot{c}_2(t)y_2(t)+\cdots+\dot{c}_n(t)y_n(t)=0\nonumber\\
	       &\cdots\cdots\cdots\nonumber\\
  \dot{c}_1(t)y^{(n-1)}_1(t)&+\dot{c}_2(t)y^{(n-1)}_2(t)+\cdots
  +\dot{c}_n(t)y^{(n-1)}_n(t)=g(t)
	       \label{eq:SystemCs}
\end{align}
Die Koeffizientendeterminante von (\ref{eq:SystemCs}) ist die
Wronskideterminante von der wir in (\ref{eq:WronskianGeneral}) gezeigt haben,
dass sie von null verschieden ist. Damit ist sichergestellt, dass
(\ref{eq:SystemCs}) eine Lösung hat.
%
\begin{example}{Variation der Konstanten}
  \[\ddot{y}(t)+y(t)=\frac{1}{\cos(t)}\]
  %TODO die homogene Lösungen sind ?
  Die allgemeine Lösung ist $y(t)=c_1\cos(t)+c_2\sin(t)$ und wir machen für die
  spezielle Lösung den Ansatz
  \[y(t)=c_1(t)\cos(t)+c_2(t)\sin(t)\]
  und erhalten das lineare Gleichungssystem
  \begin{align*}
    \dot{c}_1(t)\cos(t)+\dot{c}_2(t)\sin(t)&=0\\
    -\dot{c}_1(t)\sin(t)+\dot{c}_2(t)\cos(t)&=\frac{1}{\cos(t)}
  \end{align*}
  Die spezielle Lösung ist damit
  \[y(t)=\cos(t)\ln(|\cos(t)|)+t\sin(t)\]
\end{example}
%
\pagebreak 
\section{Lösung der inhomogenen Gleichung: Operatormethode}
Für die Ableitung benutzen wir als Abkürung eine Operatorschreibweise der Form
$\frac{dy(t)}{dt}=D[y(t)]$. Entsprechend soll $D^{-1}[y(t)]=\int y(t)dt$
verwendet werden. Wir wissen ja aus dem Haupsatz der Differential- und
Integralrechnung, dass die Ableitung die Umkehrung des Integrals ist. Für einen
Operator $D$ schreiben wir dessen Inverse als $D^{-1}$. Ein r-faches
Hintereinanderausführen eines Operators können wir also mit $D^r$ schreiben.
\begin{note}{}
  Mit dieser Schreibweise wird uns klar, dass z.B. die Exponentialfunktion der
  Ableitung
  \[e^\frac{d}{dt}=e^D=1+D+\frac{D^2}{2}+\dots=\sum_{\nu=0}^\infty\frac{D^\nu}{\nu!}\]
  bedeutet. 
  
  Merke: Funktionen von Operatoren sind über ihre Reihenentwicklung
  definiert.

  Eine offensichtliche Anwendung sehen wir mit der Taylorentwicklung von
  $y(t+\tau)$ um den Zeitpunkt $t$
  \[y(t+\tau)=y(t)+\left.\frac{\tau}{1!}\dot{y}(s)\right|_{s=t}+\left.\frac{\tau^2}{2!}\ddot{y}(s)\right|_{s=t}
	+\dots=\sum\limits_{\nu=0}^{\infty}\left.\frac{\tau^\nu}{\nu!}\frac{d^\nu}{ds^\nu}y(s)\right|_{s=t}=
	e^{\tau\frac{d}{dt}}y(t)\]
  Wir nennen $e^{\tau\frac{d}{dt}}=e^{\tau D}$ den Verschiebungs- oder
  Shiftoperator.
\end{note}
Wir betrachten die folgende Differentialgleichung erster Ordnung mit konstanten
Koeffizienten 
\begin{equation}
  \dot{y}(t)-\lambda y(t)=g(t)\mbox{ oder in Operatorform }(D-\lambda)y(t)=g(t)
  \label{eq:RelaxatorOperatorform}
\end{equation}
Wir wissen, dass (\ref{eq:RelaxatorOperatorform}) als homogene Lösung
$y(t)=e^{\lambda t}$ hat, also versuchen wir es mit dem
Produktansatz\footnote[1]{Eingesetzt ergibt dies $D[e^{\lambda
t}D^{-1}[e^{-\lambda t}g(t)]]-\lambda e^{\lambda t}D^{-1}[e^{-\lambda
t}g(t)]=g(t)$. Wenn wir die Operation $D$ ausführen - Achtung Produktregel -
und $DD^{-1}=1$ beachten, dann erhalten wir $g(t)=g(t)$}
\begin{equation}
  y(t)=e^{\lambda t}\cdot D^{-1}[e^{-\lambda t}g(t)]
  \label{eq:Produktansatz}
\end{equation}
Wir verifizieren das durch einsetzen in (\ref{eq:RelaxatorOperatorform}). Nun
wenden wir das auf die Differentialgleichung $(D-\lambda)^2[y(t)]=g(t)$ an und
machen hierfür den entsprechenden Produktansatz
\begin{equation}
  y(t)=e^{\lambda t}\cdot D^{-2}[e^{-\lambda t}g(t)]
  \label{eq:Produktansatz2Ordnung}
\end{equation}
und zeigen wieder, dass (\ref{eq:Produktansatz2Ordnung}) die
Differentialgleichung erfüllt. nun stellt sich die Frage nach der
Verallgemeinerung, d.h. was geschieht wenn die Differentialgleichung 
\begin{equation}
  \ddot{y}(t)-\alpha\dot{y}(t)+\beta y(t)=g(t)
  \label{eq:DGL2Ordnung2Koeffizienten}
\end{equation}
vorliegt, woebei wir $\alpha=\lambda_1+\lambda_2$ und $\beta=\lambda_1\lambda_2$
schreiben. Damit haben wir
\begin{equation}
  (D-\lambda_1)(D-\lambda_2)[y(t)]=g(t)
  \label{eq:DGG2KoeffizientenOperatorform}
\end{equation}
Die Gleichung (\ref{eq:DGL2Ordnung2Koeffizienten}) können wir mit $(D-\lambda_2)[y(t)]=y^*(t)$ 
in der Form $(D-\lambda_1)[y^*(t)]=g(t)$ schreiben.

Nun lösen wir das ganze sukzessive auf
\begin{align} 
  y^*(t)&=e^{\lambda_1t}\cdot D^{-1}\left[e^{-\lambda_1 t}g(t)\right]\nonumber\\
  y(t)  &=e^{\lambda_2t}\cdot D^{-1}\left[e^{-\lambda_2 t}y^*(t)\right]=\nonumber\\
        &=e^{\lambda_2t}\cdot D^{-1}\left[e^{(\lambda_1-\lambda_2) t}\cdot
      D^{-1}\left[e^{-\lambda_1 t}g(t)\right]\right]
  \label{eq:DGL2OrdnungSolu}
\end{align}
Damit ist auch der allgemeine Fall für ein Polynom beliebiger Ordnung 
\[ P(D)[y(t)]=(D-\lambda_1)^{r_1}\cdot(D-\lambda_2)^{r_2}\cdots 
(D-\lambda_n)^{r_n}[y(t)]\]
zugänglich, nämlich
\begin{equation}
  y(t)=e^{\lambda_s t}\cdot D^{-r_s}[e^{(\lambda_{s-1}-\lambda_s)t}\cdot D^{-r_{s-1}}
  [\dots e^{(\lambda_1-\lambda_2)t}D^{-r_1}[e^{-\lambda_1t}g(t)]\dots]]
  \label{eq:SoluInhomogenOperator}
\end{equation}
\begin{example}{Ganzrationale Inhomogenität}
  \[\dddot{y}(t)-3\dot{y}(t)-2y(t)=4t^2-2\]
  Hat die characteristische Gleichung
  \[(\lambda^3-3\lambda-2)=(\lambda-2)(\lambda+1)^2=0\]
  Die allgemeine Lösung der homogenen Differentialgleichung lautet somit
  \[y(t)=c_1e^{2t}+c_2e^{-t}+c_3te^{-t}\]
  Mit (\ref{eq:SoluInhomogenOperator}) erhalten wir
  \[
    y(t)=e^{-t}\cdot D^{-2}[e^{3t}\cdot D^{-^1}[e^{-2t}(4t^2-2)]]=-2t^2+6t-8
  \]
  Damit ist die allgemeine Lösung gegeben durch
  \[ y(t)=c_1e^{2t}+c_2e^{-t}+c_3te^{-t}-2t^2+6t-8 \]
  Die $c_i$ werden durch die Anfangsbedingungen festgelegt.
\end{example}
Hieraus leiten wir eine allgemeine herangehensweise bei Polynomen als
Inhomogenitäten ab.
\begin{note}{Spezielle Lösung bei speziellen Inhomogenitäten}
  \begin{enumerate}
    \item Ist die Inhomogenität gegeben als ein Polynom n-ter Ordnung, so setzt
      man die spezielle Lösung der Differentialgleichung als allgemeines
      Polynom n-ter Ordnung an und bestimmt die Polynomialkoeffizienten durch
      einsetzen in die inhomogene Differetialgleichung und
      Koeffizientenvergleich.  Wir haben 2 Fälle zu betrachten
    \item Ist $g(t)=e^{\mu t}$ und $\mu\ne\lambda_i$ dann erhalten wir durch
      einsetzen in (\ref{eq:SoluInhomogenOperator}) 
  \[ y(t)=\frac{e^{\mu t}}{(\mu-\lambda_1)^{r_1}
            \cdots(\mu-\lambda_s)^{r_s}}\]
	  \item Liegt hingegen Resonanz vor, z.B.  $\mu=\lambda_s$, dann schreiben wir (\ref{eq:SoluInhomogenOperator})
	    \[y(t)=e^{\lambda_s t}\cdot D^{-r_s}
	    \left[\frac{e^{(\mu-\lambda_s)\cdot t}}{(\mu-\lambda_1)^{r_1}\dots(\mu-\lambda_{s-1})^{r_{s-1}}}\right]\]
	    Beachte, die Reihenfolge der Operatoren $(D-\lambda_i)^{r_i}$
	    kann beliebig vertauscht werden und wir erhalten im Resonanzfall
	    \[y(t)=e^{\mu t}\cdot D^{-r_s}
	      \left[\frac{1}{(\mu-\lambda_1)^{r_1}\dots(\mu-\lambda_{s-1})^{r_{s-1}}}\right]\]
	    Was sich wiederum integrieren lässt zu
	\[ y(t)=\frac{t^{r_s}e^{\mu t}}{(\mu-\lambda_1)^{r_1} \cdots(\mu-\lambda_{s-1})^{r_{s-1}}}
	\cdot\frac{1}{r_s!}\]
  \end{enumerate}
\end{note}
Wir fassen die allgemeine Vorgehensweise im Diagramm in Abbildung
\ref{fig:FlowchartSolu} zusammen.
%
\begin{figure}[htb]
  \centering
  \tikzstyle{decision} = [diamond, draw, 
    text width=4.5em, text badly centered, node distance=3cm, inner sep=0pt]
  \tikzstyle{block} = [rectangle, draw, 
    text width=8em, text centered, rounded corners, minimum height=4em]
  \tikzstyle{line} = [draw, -latex']
  \tikzstyle{cloud} = [draw, ellipse, node distance=3cm,
    minimum height=2em]
\begin{tikzpicture}[node distance = 2cm, auto]
    % Place nodes
    \node [block] (init) {DGL-Löser};
    \node [block, left of=init, node distance=4cm] (dgl) {Differential- gleichung};
    \node [block, right of=init, node distance=4cm] (ics) {Anfangsbe- dingungen};
    \node [decision, below of=init] (hominhom) {ist es eine homogene DGL?};
    \node [block, right of=hominhom, node distance=4cm] (partsolu) {Suche spezielle L\"osung};
    \node [block, below of=hominhom, node distance=3cm] (homsolu) {Suche homogene L\"osung(en)};
    \node [block, below of=homsolu, node distance=2cm] (hometinhom) {Addiere
           homogene und inhomogene L\"osung(en)};
    \node [block, below of=hometinhom, node distance=2cm] (stop) {Erf\"ulle Randbedingungen};
         % Draw edges
    \path [line] (init) -- (hominhom);
    \path [line] (hominhom) -- node {immer} (homsolu);
    \path [line] (hominhom) -- node {yes} (partsolu);
    \path [line] (homsolu) -- (hometinhom);
    \path [line] (hometinhom) -- (stop);
    \path [line] (dgl) -- (init);
    \path [line] (ics) -- (init);
    \path [line] (partsolu) |- (stop);
\end{tikzpicture}
\caption{Flussdiagramm zur Lösung inhomogener linearer Differentialgleichungen.}
  \label{fig:FlowchartSolu}
\end{figure}
Zur Erfüllung der Randbedingungen gibt es zwei Vorgehensweisen, die wir im
folgenden diskutieren wollen.

In der ersten Vorgehensweise überlagern wir die gefundene spezielle Lösung mit
der homogenen Lösung, die ja noch zu wählende Konstanten beinhaltet
$y(t)=y_{hom}(t)+y_{par}(t)$. Dann setzen wir diese Lösung in die
Randbedingungen ein und bestimmen so die Konstanten. Dies ist die sichere
Variante.

Für die zweite Vorgehensweise erinnern wir uns an die Variation der Konstanten.
Dort hatten wir für eine inhomogene Gleichung erster Ordnung
$(D-\lambda)y(t)=f(t)$ den Ansatz $y_{par}(t)=c(t)\cdot e^{\lambda t}$ gemacht.
Die Besitummungsgleichung für $c(t)$ lautete $\dot c(t)=e^{\lambda t}\cdot
f(t)$. Nun integrieren wir über t in den Grenzen von $0$ bis $t$. Damit
erhalten wir $c(t)-c(0)=\int_0^te^{\lambda t'}\cdot f(t')dt'$. Die Konstante
$c(0)$ ist unbestimmt. Wir wählen $c(0)=0$. Die Integration in den Grenzen
wenden wir auch auf (\ref{eq:SoluInhomogenOperator}) an. Damit wird
(\ref{eq:SoluInhomogenOperator}) zu
\begin{equation}
  y_p(t)=e^{\lambda_n t}
\int\limits_0^t e^{(\lambda_{n-1}-\lambda_n) t'}
\int\limits_0^{t'} e^{(\lambda_{n-2}-\lambda_{n-1}) t''}\dots
\int\limits_0^{r} e^{(\lambda_{1}-\lambda_{2}p) p}\
\int\limits_0^{p} e^{-\lambda_{1}q}f(q)dqdp\dots dt''dt'
  \label{eq:SoluInhomogenOperator0}
\end{equation}
Wie man leicht sieht, gilt $y_p(0)=0$, $\dot y_p(0)=0$,\ldots,
$y^{(n-1)}_p(0)=0$ (wir zeigen dies in der Anmerkung unten). Damit bleibt die
Aufgabe, die $c_i$ der homogenen Lösung so zu bestimmen, dass die homogene
Lösung die Anfangsbedingungen erfüllt.
\begin{note}{Die Ableitungen von $y_p(t)$ aus (\ref{eq:SoluInhomogenOperator0}) bei $t=0$}
Wir leiten (\ref{eq:SoluInhomogenOperator0}) einmal nach $t$ ab, dann erhalten
wir (Kettenregel beachten)
\[ \dot y_p(t)=\lambda_n y_p(t)+
e^{\lambda_n t}\frac{d}{dt}
\int\limits_0^t e^{(\lambda_{n-1}-\lambda_n) t'}
\int\limits_0^{t'} e^{(\lambda_{n-2}-\lambda_{n-1}) t''}\dots
\int\limits_0^{r} e^{(\lambda_{1}-\lambda_{2}p) p}\
\int\limits_0^{p} e^{-\lambda_{1}q}f(q)dqdp\dots dt''dt'\
\]
Unter dem Integral über $dt'$ steht eine Funktion $g(t')$, n'est-ce pas? Nämlich
\[g(t')=e^{(\lambda_{n-1}-\lambda_n) t'}
\int\limits_0^{t'} e^{(\lambda_{n-2}-\lambda_{n-1}) t''}\dots
\int\limits_0^{r} e^{(\lambda_{1}-\lambda_{2}p) p}\
\int\limits_0^{p} e^{-\lambda_{1}q}f(q)dqdp\dots dt''
\]
Also haben wir
\[ \dot y_p(t)=\lambda_n y_p(t)+e^{\lambda_n t}\frac{d}{dt}
\int\limits_0^t g(t')dt'
\]
Davon wissen wir aber, dass gilt
\[ \lambda_n y_p(t)+e^{\lambda_n t} g(t)
\]
Es ist aber $y_p(0)=0$ und, wie wir leicht sehen, auch $g(0)=0$. Der geneigte Leser möge letzte Aussage beweisen. Also ist auch $\dot y_p(0)=0$. Das geht weiter bis zur Ableitung der Ordnung $n-1$.
\end{note}
\chapter{Systeme linearer gewöhnlicher Differentialgleichungen}
\Comment{4. Vorlesung}
Wir benutzen die Operatorschreibweise, um ein allgemeines System von
Differentialgleichungen beliebiger Ordnung anzugeben. Wir betrachten somit
Systeme der Art
\begin{equation}
  P_{i_1}(D)[y_1(t)]+P_{i_2}(D)[y_2(t)]+\dots+P_{i_n}(D)[y_n(t)]=f_i(t)\hfill
  (i=1,2,\dots,m)
  \label{eq:SystemGeneral}
\end{equation}
Um (\ref{eq:SystemCs}) zu lösen gibt es zwei Strategien: Erstens könnten wir
aus dem vorliegenden System eine Differentialgleichung höherer Ordnung machen
oder zweitens könnten wir versuchen auf ein System von gekoppelten
Differentialgleichungen erster Ordnung überzugehen.

Anhand zweier Beispiele wollen wir die Vor- und Nachteile beider Methoden
untersuchen und aus den Nachteilen gegebenenfalls Rückschlüsse auf die Modelle
ziehen.
\begin{example}{Rückführung auf DGL höherer Ordnung}
  Ähnlich dem Gauß'schen Eliminationsverfahren schreiben wir
  \begin{equation*}
    \begin{array}{rrl|l}
      D^2[y_1(t)]+&y_2(t)&=0&D^2\\
      y_1(t)+&D^2[y_2(t)]&=0&\cdot(-1)
    \end{array}
  \end{equation*}
  und erhalten die DGL $D^4[y_1(t)]-y_1(t)=0$.

Versuchen wir die Methode hingegen auf das System
\begin{align*}
  (D+1)[y_1(t)]+(D^2-1)[y_2(t)]&=0\\
  y_1(t)+(D-1)[y_2(t)]&=0
\end{align*}
anzuwenden, gelingt die Elimination von $y_1(t)$ oder $y_2(t)$ nicht, denn die
beiden Gleichungen sind linear abhängig (warum?).
\end{example}
Wie wir aus dem Gegenbeispiel gesehen haben, ist die Tatsache der linearen
Abhängigkeit auf eine nicht eindeutige Formulierung der Problemstellung
zurückzuführen. Denn entweder ist $y_1(t)$ oder $y_2(t)$ frei wählbar und die
Gleichung dann nach der jeweils anderen Funktion zu lösen. Diese Wahlfreiheit
bedeutet aber einen Mangel an Information über das Modellsystem. Für eine
eindeutige Lösung müssen wir uns Gedanken über weitere Bedingungen machen, die
wir dem System auferlegen, damit es eindeutig mit Hilfe der
Differentialgleichungen formuliert werden kann. Das beudeutet dann noch nicht,
dass es lösbar sein wird.
\begin{example}{Rückführung auf ein System linearer DGL 1. Ordnung}
  Aus dem System
  \begin{align*}
    \ddot{y}_1(t)+y_2(t)=0\\
    \ddot{y}_2(t)+y_1(t)=0
  \end{align*}
  erhalten wir mit $\dot{y}_1(t)=y_3(t)$ und $\dot{y}_2(t)=y_4(t)$ das System
  \begin{align*}
    \dot{y}_1(t)=y_3(t)\\
    \dot{y}_2(t)=y_4(t)\\
    \dot{y}_3(t)=-y_2(t)\\
    \dot{y}_4(t)=-y_1(t)
  \end{align*}
  Versuchen wir allerdings diese Methode auf das System
    \begin{align*}
      \ddot{y}_2(t)+\dot{y}_1(t)+y_1(t)-y_2(t)=0\\
      \dot{y}_2(t)+y_1(t)-y_2(t)=0
  \end{align*}
  anzuwenden, dann lässt sich dieses nicht in ein System der gewünschten Art
  überführen. In diesem Fall lässt sich die Funktion $y_1(t)$ frei wählen und
  dann das lineare System lösen. Da aber von vornherein $y_1(t)$ nicht
  eindeutig bestimmt ist, gilt das oben gesagte.
\end{example}
Wir werden uns im folgenden daher nur mit der Form
\[
  \dot{y}_i(t)=a_{i1}y_1(t)+a_{i2}y_2(t)+\dots+a_{in}y_n(t)\hfill (i=1,2,\dots,m)
\]
befassen.
\section{Homogene Systeme erster Ordnung}
Ein homogenes System von linearen Differentialgleichungen erster Ordnung hat
die Form
\begin{equation}
  \begin{pmatrix}
    \dot{y}_1(t)\\
    \dot{y}_2(t)\\
    \dots\\
    \dot{y}_n(t)
  \end{pmatrix}
  =
  \begin{pmatrix}
    a_{11}&a_{12}&\dots &a_{1n}\\
    a_{21}&a_{22}&\dots &a_{2n}\\
    \multicolumn{4}{c}\dotfill\\
    a_{n1}&a_{n2}&\dots &a_{nn}\\
  \end{pmatrix}
  \begin{pmatrix}
    y_1(t)\\
    y_2(t)\\
    \dots\\
    y_n(t)
  \end{pmatrix}
  \label{eq:System1stOrder}
\end{equation}
Das System (\ref{eq:System1stOrder}) lässt sich mit Hilfe der
Operatorschreibweise auch so darstellen
\begin{equation}
  \begin{pmatrix}
    a_{11}-D&a_{12}&\dots &a_{1n}\\
    a_{21}&a_{22}-D&\dots &a_{2n}\\
    \multicolumn{4}{c}\dotfill\\
    a_{n1}&a_{n2}&\dots &a_{nn}-D\\
  \end{pmatrix}
  \begin{pmatrix}
    y_1(t)\\
    y_2(t)\\
    \dots\\
    y_n(t)
  \end{pmatrix}
  =0
  \label{eq:System1stOperator}
\end{equation}
Oder wir schreiben es ganz einfach als
\begin{equation}
  \dot{\vec{y}}(t)=A\cdot\vec{y}(t)
  \label{eq:SystemVectorMatrixNotation}
\end{equation}
Es gibt genau eine Lösung $(y_1(t),y_2(t),\dots,y_n(t))$ von
(\ref{eq:System1stOrder}), so dass die zu einem beliebigen Zeitpunkt $t_0$
gewählten Bedingungen
\begin{equation}
  y_1(t_0)=b1,\:y_2(t_0)=b_2,\dots,y_n(t_0)=b_n
  \label{eq:ICsSystem}
\end{equation}
erfüllt sind. Wir werden den Beweis hier nicht antreten und verweisen
stattdessen auf die einschlägigen Lehrbücher.

Für die Lösung der Systeme (\ref{eq:System1stOrder}) benutzen wir den
Exponentialansatz $e^{\lambda t}$. Da wir $n$ Funktionen zu bestimmen haben,
schreiben wir diesen gleich in der Form
\begin{equation}
  y_i(t)=c_ie^{\lambda t}\hfill (i=1,2,\dots,n\mbox{ und }c_i\in\mathbb{R})
  \label{eq:SystemAnsatz}
\end{equation}
Wir setzen den Ansatz in (\ref{eq:System1stOrder}) ein und erhalten das
Gleichungssystem
\begin{align}
    (a_{11}-\lambda)c_1+a_{12}c_2+\dots+a_{1n}c_n&=0\nonumber\\
    a_{21}c_1+(a_{22}-\lambda)c_2+\dots+a_{2n}c_n&=0\nonumber\\
    \dots\dots\dots\dots\nonumber\\
    a_{21}c_1+a_{22}c_2+\dots+(a_{2n}-\lambda)c_n&=0
  \label{eq:System4cs}
\end{align}
Oder einfacher
\begin{equation}
  (A-\lambda\mathbbm{1})\cdot\vec{c}=0
  \label{eq:Vektorform}
\end{equation}
Aus der linearen Algebra wissen wir, dass (\ref{eq:Vektorform}) nur eine
Lösung hat, wenn gilt $|A-\lambda\mathbbm{1}|=0$ und die für die Lösungen
gewählten Werte der $\lambda$ sind die Eigenwerte der Koeffizientenmatrix $A$.

Um (\ref{eq:Vektorform}) zu lösen erinnern wir uns an Eigenwerte, Eigenvektoren
und Normalformen quadratischer Matrizen. Wir wissen es gibt drei Fälle:
\begin{enumerate}
  \item Alle Eigenwerte $\lambda_i$ sind voneinander verschieden.
  \item Es treten mehrfache Eigenwerte auf. Der Rang der Matrix
    $A-\lambda\mathbbm{1}$ fällt für jeden mehrfachen Eigenwert um soviel
    gegenüber $n$, wie die Vielfachheit des entsprechenden Eigenwertes beträgt.
  \item Es tritt mindestens ein mehrfacher Eigenwert auf für den der Rangabfall
    nicht die obere Regel erfüllt.
\end{enumerate}
Die ersten beiden Fälle erlauben uns die Matrix $A$ auf Diagonalform 
\begin{equation}
  \Lambda=
  \begin{pmatrix}
    \lambda_1&0&\dots&0\\
    0&\lambda_2&\dots&0\\
    \multicolumn{4}{c}\dotfill\\
    0&0&\dots&\lambda_n
  \end{pmatrix}
  \label{eq:Diagonalform}
\end{equation}
zu bringen. Was bedeutet das für unser Gleichungssystem? Wir schreiben
$\vec{y}^T$ als Transformierte eines Vektors $\vec{z}^T$ dergestalt
$\vec{y}^T=C\cdot\vec{z}^T$ und damit auch $\dot{\vec{y}}^T = C\cdot
\dot{\vec{z}}^T$. Damit erhalten wir 
\[ C\cdot \dot{\vec{z}}^T = A\cdot C\cdot\vec{z}^T \]
und wenn wir mit $C^{-1}$ multiplizieren ergibt sich
\begin{equation}
  \dot{\vec{z}}^T = C^{-1}A\cdot C\cdot\vec{z}^T
  \label{eq:Diagonalsystem}
\end{equation}
Der Vorteil erschließt sich uns sofort: wir lösen das einfache Diagonalsysteme
(\ref{eq:Diagonalsystem}) und transformieren dann zurück um $\vec{y}(t)$ zu
erhalten.
%

\textbf{Fall 1: }Wir bestimmen zunächst die Eigenvektoren $(c_{i1}, c_{i2},
\dots, c_{in})$, die zu den jeweiligen Eigenwerten $\lambda_i$ gehören. Diese
bilden die Spalten der Transformationsmatrix 
\begin{equation}
  C=\begin{pmatrix}
    c_{11}&c_{21}&\dots&c_{n1}\\
    c_{12}&c_{22}&\dots&c_{n2}\\
    &&\dots&\\
    c_{1n}&c_{2n}&\dots&c_{nn}
  \end{pmatrix}
  \label{eq:Tranformationsmatrix}
\end{equation}
mit deren Hilfe wir die Matrix $A$ in die Diagonalform $\Lambda$ übeführen können
\begin{equation}
  C^{-1}\cdot A\cdot C=\Lambda
%  \label{eq:Diagonalform}
\end{equation}
Mit (\ref{eq:Tranformationsmatrix}) können wir die allgemeine Lösung angeben
\begin{equation}
  \begin{pmatrix}y_1(t)\\y_2(t)\\\dots\\y_n(t)\end{pmatrix}=
  \begin{pmatrix}
    c_{11}d_1&c_{21}d_2&\dots&c_{n1}d_n\\
    c_{12}d_1&c_{22}d_2&\dots&c_{n2}d_n\\
    &&\dots&\\
    c_{1n}d_1&c_{2n}d_2&\dots&c_{nn}d_n\\
  \end{pmatrix}\cdot
  \begin{pmatrix}e^{\lambda_1t}\\e^{\lambda_2t}\\\dots\\e^{\lambda_nt}\end{pmatrix}
  \label{eq:LsgAllgFall1}
\end{equation}
Die beliebig wählbaren reellen Konstanten $d_i$ benutzen wir zur Erfüllung der
Anfangsbedingungen.

\textbf{Fall 2: }Es treten mehrfache Eigenwerte auf und der Rangabfall der
Koeffizientenmatrix des Gleichungssystems aus (\ref{eq:System4cs}) ist gleich
dem Vielfachen des Eigenwertes. Auch in diesem Fall lässt sich $A$ in die
Diagonalform bringen. Wir erhalten parametrische Lösungen für die
Eigenvektroren der mehrfachen Eigenwerte, die wir so wählen können, dass sie
linear unabhängig sind. Das weitere Vorgehen ist genauso, wie in Fall 1.

\textbf{Fall 3: }In diesem Fall können wir die Matrix $A$ nicht auf
Diagonalform jedoch auf die Jordansche Normalform bringen.
\begin{equation*}
  J=\begin{pmatrix}
    \lambda_1&\delta_1&0&\dots&0&0\\
    0&\lambda_2&\delta_2&\dots&0&0\\
    0&0&\lambda_3&\dots&0&0\\
    &&&\dots&&\\
    0&0&0&\dots&\lambda_{n-1}&\delta_{n-1}\\
    0&0&0&\dots&0&\lambda_n\\
  \end{pmatrix}
\end{equation*}
Die $\lambda_i$ sind die Eigenwerte der Matrix $A$ und die $\delta_i$ entweder
$0$ oder $1$. Es gibt eine Transformationsmatrix $C$, die $A$ auf die
Jordanform gemäß \[ C^{-1}\cdot A\cdot C=J \] abbildet. Für einen vielfachen
Eigenwert $\lambda_k$ mit der Multiplizität $k$ bilden wir aus den Lösungen der
Gleichungen $(A - \lambda_k \mathbbm{1})^2 \vec{c}^T = 0$, $(A - \lambda_k
\mathbbm{1})^3\vec{c}^T = 0$, \dots, $(A - \lambda_k \mathbbm{1})^k \vec{c}^T =
0$ linear unabhängige Vektoren, die wir zur Transformationsmatrix
zusammenbauen. Wir bringen das System (\ref{eq:System1stOrder}) auf die Jordanform und lösen dann die Gleichung
\begin{equation}
  \begin{pmatrix}\dot{z}_1\\\dot{z}_2\\\dot{z}_3\\\dots\\\dot{z}_{n-1}\\\dot{z}_{n}\end{pmatrix}=
  \begin{pmatrix}
    \lambda_1&\delta_1&0&\dots&0&0\\
    0&\lambda_2&\delta_2&\dots&0&0\\
    0&0&\lambda_3&\dots&0&0\\
    &&&\dots&&\\
    0&0&0&\dots&\lambda_{n-1}&\delta_{n-1}\\
    0&0&0&\dots&0&\lambda_n\\
  \end{pmatrix}
  \begin{pmatrix}z_1\\z_2\\z_3\\\dots\\z_{n-1}\\z_{n}\end{pmatrix}
  \label{eq:JordanNormalForm} 
\end{equation}
Wieder lösen wir das System (\ref{eq:JordanNormalForm}) nach $\vec{z}^T$ auf
und transformieren danach wieder zurück\newline $\vec{y}^T=C\cdot\vec{z}^T$. 
\section{Inhomogene Systeme erster Ordnung}
In diesem Abschnitt beschäftigen wir uns mit den inhomogenen Systemen von
Differentialgleichungen erster Ordnung der Form
\begin{equation}
  \dot{\vec{\mathbf{x}}}=\mathbf{A}\vec{\mathbf{x}}+\vec{\mathbf{b}}
  \label{eq:InhomSystem}
\end{equation}
Eine Fundamentallösung von (\ref{eq:InhomSystem}) in der Form einer
Matrixfunktion $\boldsymbol{\Phi}(t)$ gehorcht der Gleichung 
\begin{equation}
  \dot{\boldsymbol{\Phi}}(t)=\mathbf{A}\boldsymbol{\Phi}(t)
  \label{eq:Fundamentalmatrix}
\end{equation}
$\boldsymbol{\Phi}(t)=e^{At}$ ist eine Fundamentalmatrix, welche die Bedingung
$\boldsymbol{\Phi}(0)=\mathbbm{1}$ erfüllt. Für (\ref{eq:InhomSystem}) haben wir
$\boldsymbol{\Phi}(t)=e^{At}\mathbf{C}$, mit einer beliebigen nichtsingulären
Matrix $\mathbf{C}$.  Wenn aber mit $\boldsymbol{\Phi(t)}$ eine Fundamentalmatrix
gefunden ist, dann ist es leicht ein System wie (\ref{eq:InhomSystem}) mit den
Anfangsbedingungen $\mathbf{x}(0)=\mathbf{x}_0$ zu lösen
\begin{equation}
  \mathbf{x}(t)=\boldsymbol{\Phi}(t)\boldsymbol{\Phi}^{-1}(0)\mathbf{x}_0
  +\int\limits_0^t\boldsymbol{\Phi}(t)\boldsymbol{\Phi}^{-1}(\tau)\vec{\mathbf{b}}(\tau)d\tau
  \label{eq:LoesInhom}
\end{equation}
Setzen wir hier $\boldsymbol{\Phi}(t)=e^{\mathbf{A}t}$ ein, dann erhalten wir
\[ \mathbf{x}(t)=e^{\mathbf{A}t}\mathbf{x}_0
   +e^{\mathbf{A}t}\int\limits_0^te^{-\mathbf{A}\tau}\;\vec{\mathbf{b}}(\tau)d\tau\]

   \begin{example}{Getriebener harmonischer Oszillator}
     \[\ddot{x}(t)+x(t)=f(t)\]
     kann als inhomogenes System erster Ordnung geschrieben werden. Wir setzen
     $y_1(t)=x(t)$ und $y_2(t)=\dot{x}(t)$
     \begin{align*}
       \dot{y}_1(t)&=y_2(t)\\
       \dot{y}_2(t)&=-y_1(t)+f(t)
     \end{align*}
     In diesem Fall ist
     \[\mathbf{A}=
       \begin{pmatrix}
	 0&1\\ -1&0
       \end{pmatrix}\quad
       \mbox{ und }\quad
       \vec{\mathbf{b}}(t)=\begin{pmatrix}0\\ f(t)\end{pmatrix}
     \]
     Wir wissen, dass
     \[e^{\mathbf{A}t}=
       \begin{pmatrix}
	 \cos(t)&-\sin(t)\\ \sin(t)&\cos(t)
       \end{pmatrix}
       =R(t)
     \]
     eine Rotationsmatrix ist und es ist
     \[e^{-\mathbf{A}t}=
       \begin{pmatrix}
         \cos(t)&\sin(t)\\-\sin(t)&\cos(t)
       \end{pmatrix}
       =R(-t)
     \]
     Wir erhalten als Lösung
     \[\mathbf{y}(t)=R(t)\mathbf{y}_0 + 
       R(t)\int\limits_0^t\begin{pmatrix}f(\tau)\sin(\tau)\\f(\tau)\cos(\tau)\end{pmatrix}d\tau
     \]
     Das ursprüngliche Oszillatorsystem hat also die Lösung $x(t)=y_1(t)$
     \[ x(t)=x_0\cos(t)-v_0\sin(t)+\int\limits_0^tf(\tau)\sin(\tau -t)d\tau \]
   \end{example}
\input{chapters/5_Integraltrafo}
\chapter{Theorie der linearen Antwort}
\Comment{6. Vorlesung}
Dieses äusserst wichtige Thema behandelt die linearen passive Systeme, wie wir
sie aus der Elektrotechnik kennen. Hierbei bedeutet ``Passivität'', dass das
System im zeitlichen Mittel nur Arbeit bzw. Energie aufnimmt, wenn äussere
Kräfte darauf einwirken. Dieser Abschnitt folgt im wesentlichen der
Darstellung, wie sie in (Jaeckle, 1978) zu finden ist.
\section{Der gedämpfte harmonische Oszillator: 
\it Exkursion zur Drosophila melanogaster der Physik}\label{sec:gho}
Da der harmonische Oszillator in zahlreichen Variationen und auf den
verschiedensten Gebieten immer wieder anzutreffen ist, wollen wir ihm hier
einen extra Unterabschnitt widmen. 

Die Auslenkung $x(t)$ eines gedämpften harmonischen Oszillators als Antwort
auf eine von außen angelegte zeitabhängige Kraft $f(t)$ ist gegeben durch
die Differentialgleichung
\begin{equation}
  m\left(\frac{d^2}{dt^2}+\beta\frac{d}{dt}+\omega_0^2\right)x(t)=f(t)  
  \label{eq:gdho}
\end{equation}
wobei $\beta$ die Dämpfungskonstante bezeichne und $\omega_0/2\pi$ die
Eigenfrequenz des Oszillators ohne Dämpfung.

Für eine beliebige äussere Kraft $f(t)$ existiert eine sogenannte {\it
lineare Antwortfunktion} oder auch {\it Greensche Funktion} $G(\tau)$ genannt,
mit Hilfe derer wir die Lösung des Anfangswertproblems mit
$\lim_{t_0\rightarrow 0}x(t_0)=0$ und $\lim_{t_0\rightarrow 0}\dot{x}(t_0)=0$
auf folgende Weise schreiben können
\begin{equation}
  x(t)=\int_{-\infty}^{\infty}G(t-t')f(t')dt'
  \label{eq:solugreengdho}
\end{equation}
Die Greensche Funktion ist eine lineare Abbildung der Funktion $f(t)$ auf die
Funktion $x(t)$ (man zeige dass (\ref{eq:solugreengdho}) tatsächlich eine
lineare Abbildung ist!).

Für $\frac{\beta}{2}<\omega_0$, dem Fall schwacher Dämpfung, lautet die
Greensche Funktion 
\begin{equation}
  G(\tau)=\frac{1}{m\omega_r}e^{-\frac{\beta}{2}\tau}\sin(\omega_r\tau)\Theta(\tau)
  \label{eq:Greenfunction}
\end{equation}
wobei $\omega_r=\sqrt{\omega_0^2-\frac{\beta^2}{4}}$ die Resonanzfrequenz im
Dämpfungsfall bezeichne und 
\begin{equation*}
  \Theta(\tau)=\left\{\begin{array}{cc}1&\mbox{ für }\tau>0\\0&\mbox{ sonst. }\end{array}\right.
\end{equation*}
Diese Lösung des harmonischen Oszillators mit Hilfe einer linearen
Antwortfunktion oder der Greenschen Funktion kann in gleicher Weise auf eine
gro{\ss}e Klasse von Systemen linearer Differentialgleichungen verallgemeinert
werden.

Wir sehen im Fall des harmonischen Oszillators, dass die Greensche Funktion die
Differentialgleichung
\begin{equation}
  m\left(\frac{d^2}{d\tau^2}+\beta\frac{d}{d\tau}+\omega_0^2\right)G(\tau)=\delta(\tau)  
  \label{eq:dglGreensfunction}
\end{equation}
erfüllt. Dies zeigt man durch einsetzen von (\ref{eq:solugreengdho}) in (\ref{eq:gdho}).
Die Greensche Funktion beschreibt deshalb die Lösung der harmonischen
Oszillatorgleichung als Antwort auf einen Deltapuls, also die
Impulsantwortfunktion.
\begin{exercise}{Die Antwort des Oszillators}
Zeige mit Hilfe der Laplacetransformation, dass (\ref{eq:Greenfunction}) die
Impulsantwortfunktion des gedämpften harmonischen Oszillators ist.
\end{exercise}
Es interessiert uns besonders die harmonische Antwort des Oszillators.  Wir
suchen die Lösung von (\ref{eq:gdho}) bei einer Inhomogenität der Form
\begin{equation}\label{eq:harminh}
f(t)=\Re\{f_0\cdot e^{-izt}\}=|f_0| \cos(\omega t-\delta)e^{\eta t}\mbox{ mit }z=\omega
+i\eta\mbox{ und }\eta >0
\end{equation}
\begin{note}{Adiabatisches Einschalten}
$\eta >0$ entspricht einem langsamen ''adiabatischen'' Einschalten der externen
treibenden Kraft. Wir betrachten immer den Limes $\eta\rightarrow 0$
interessiert.
\end{note}
Wir suchen eine partikuläre Lösung der Form
\begin{equation}\label{eq:partsol}
x(t)=\Re\{x_0\cdot e^{-izt}\}=|x_0| \cos(\omega t-\delta-\phi)e^{\eta t}
\end{equation}
Wir setzen nun (\ref{eq:harminh}, \ref{eq:partsol}) in (\ref{eq:gdho}) ein und
erhalten damit für das Verhältnis der komplexen Amplituden
\begin{equation}\label{eq:dynsusz}
\frac{x_0}{f_0}=\frac{1/m}{-z^2-i\beta z+\omega_0^2}=\chi(z)
\end{equation}
Wir nennen $\chi(z)$ die {\it dynamische Suszeptibilität} des Systems, in unserem
Fall des harmonischen Oszillators. Die dynamische Suszeptibilät bestimmt die
Antwort $x(t)$ des gedämpften harmonischen Oszillators auf eine harmonische
treibende Kraft $f(t)$, die mit der Frequenz $\omega$ oszilliert. Wir können
die dynamische Suszeptibilität auch mit der Greenschen Funktion nach
(\ref{eq:solugreengdho}) erhalten. Der Vergleich der Ergebnisse zeigt, dass die
dynamische Suszeptibilität die Laplacetransformierte der Greenschen Funktion
ist
\begin{equation}\label{eq:suszept}
\chi(z) = \int_0^\infty G(t)e^{izt}dt
\end{equation}
\begin{note}{}
Zeige, dass die Form der Laplacetransformation in (\ref{eq:suszept}) äquivalent
derjenigen in (\ref{eq:Laplacetrafo}) ist!
\end{note}
Nun definieren wir uns den Fall der Resonanz und zwar auf zwei verschiedene
Arten. Einmal als die Frequenz, bei der der Oszillator mit der maximalen
Amplitude ''antwortet'' und zum anderen als diejenige, bei der seine
Leistungsaufnahme maximal ist. Für ersteres gilt
\begin{equation}\label{eq:ampres}
\lim_{\eta\rightarrow 0}|\chi(\omega +i\eta)|=
\frac{1}{m\sqrt{(\omega^2-\omega_0^2)^2+(\beta\omega)^2}}
\end{equation}
und das Maximum der Auslenkung liegt bei
$\omega_{max}=\sqrt{\omega_0^2-\beta^2/2}$, also etwas unterhalb der
natürlichen Frequenz $\omega_0$. Die momentane Leistungsaufnahme eines Systems,
die an ihm verrichtet wird, ist gegeben durch die äussere Kraft multipliziert
mit der zeitlichen Veränderung der Auslenkung - also der Geschwindigkeit - was
sich schreiben lässt als $P(t)=f(t)\cdot\dot{x}(t)$. Wir berechnen die über
eine Schwingungsperiode $T$ gemittelte Leistung als
\begin{equation}\label{eq:PMittel}
\overline{P(t)}=\frac{1}{T}\int_0^T f(t)\dot{x}(t)dt
\end{equation}
Um $\dot{x}(t)$ zu berechnen benutzen wir (\ref{eq:dynsusz}) und erhalten
zunächst 
\[x(t)=\Re [f_0\cdot\chi(z)\cdot e^{-izt}]\] 
Wir bilden die Ableitung und machen
den Gernzübergang $\eta\rightarrow 0$ und erhalten 
\begin{equation}\label{eq:xdot}
\dot{x}(t)=\omega\cdot |f_0|(-\chi'(\omega)\sin(\omega t)+\chi''(\omega)\cos(\omega t)
\end{equation}
Dabei haben wir folgenden Grenzübergang gemacht
\[ 
\lim_{\eta\rightarrow 0}\chi(\omega+i\eta)=\chi'(\omega)+i\chi''(\omega)
\]
Wenn wir dies in (\ref{eq:PMittel}) einsetzen, dann erhalten wir 
\begin{equation}\label{eq:PMittelChi}
\overline{P(t)}=\frac{1}{2}|f_0|^2\omega\chi''(\omega)
\end{equation}
Der wesentliche Faktor ist also
\begin{equation}\label{eq:PMittelFactor}
\omega\chi''(\omega)=\frac{\beta\omega^2/m}{(\omega^2-\omega_0^2)^2+\beta^2\omega^2}
\end{equation}
Die beiden Resonanzkurven für Amplituden- und Leistungsresonanz zeigen wir in
einem Diagramm in Bild \ref{fig:resonances}
\begin{figure}[!h]
\begin{center}
%\includegraphics[height=0.6\textwidth]{/Users/greiner/TeX/LectureNotes/Basics/SimulationMST/SIMOnlyD/Figures/resonanz.pdf}
\includegraphics[height=0.6\textwidth]{fig/resonanz.pdf}
\caption{Amplituden- und Leistungsresonanz eines gedämpften harmonischen Oszillators.\label{fig:resonances}}
\end{center}
\end{figure}\newpage
\begin{exercise}{Amplituden- und Leistungsresonanz}
Interpretiere die beiden Resonanzkurven aus Abbildung \ref{fig:resonances}!
\end{exercise}
%\begin{exercise}{Die Antworten des Relaxators}
%Gegeben die Differentialgleichung eines gedämpften harmonsichen Oszillators,
%wie wir sie in Abschnitt \ref{sec:gho} bereits kennen gelernt haben,
%\begin{equation}\label{eq:relaxator}
%\dot{x}(t)+\frac{1}{\tau} x(t)=f(t)
%\end{equation}
%\begin{enumerate}
%  \item Iterative Lösung: integrieren Sie (\ref{eq:relaxator}) für $f(t)=0$ und
%	$x(0)=x_0$ von $t=0$ bis $t$. Damit erhalten Sie
%	$x(t)=x_0+\int_0^t\left(-\frac{1}{\tau} x(\xi)\right)d\xi$. Wenden Sie nun den in Abschnitt
%	\ref{sec:analyticsolu} an und ermitteln Sie die Lösung.
%  \item Berechnen Sie die Antwortfunktionen für:
%	\begin{itemize}
%	  \item $f(t)=\delta(t)$ die 
%	  \item $f(t)=\Theta(t-t_0)$, wobei $\Theta(z)$ die Heavisidesche
%                Sprungfunktion darstellt, die $0$ ist für$z\le 0$ und $1$ sonst.
%	  \item $f(t)=\alpha t$, the famous "ramp response".
%	  \item $f(t)=\sin(\omega t)$, also die harmonische Antwort.
%	\end{itemize}
%  \item Das ganze mit Hilfe der Laplacetransformation. Geben Sie auch die
%	Lösung im Frequenzbereich an.
%\end{enumerate}
%%\input{Exercises/Chapter2/Relaxator.tex}
%\end{exercise}
%
\section{Die Funktion der linearen Antwort}
Im allgemeinen Fall eines linearen zeitinvarianten Systems (LTI-System) wissen
wir, dass die Antwortfunktion nur von der Zeitdifferenz $t-t'$ abhängt, d.h.
die allgemeine Antwort reduziert sich auf die Form wie beim harmonischen
Oszillator
\begin{equation}
  x(t)=\int_{-\infty}^{\infty}G_{LTI}(t,t')f(t')dt'=\int_{-\infty}^{\infty}G(t-t')f(t')dt'
  \label{eq:soluLTI}
\end{equation} 
Wir interessieren uns also nur für die Laplacetransformierte von $G(\tau)$,
wobei wir wissen, dass $G(\tau)=0$ für $\tau<0$. Dann ist die Antwort auf die
angelegte Kraft 
\begin{equation}
  x(t)=\int\limits_{-\infty}^{\infty}G(t-t')F(t')dt'
  \label{eq:Response}
\end{equation} 
$G(\tau)$ ist die reelle Systemfunktion. Wir nennen Sie die Funktion der
linearen Antwort. Das System solle sich kausal verhalten, was bedeutet, dass
$x(t)=0$ für alle Zeiten $t$, die vor dem Eischalten der Kraft $f(t)$ liegen.
Das bedeuet, dass $G(\tau)=0$ für $\tau<0$.

Ausserdem sollen $x(t)$ und $f(t)$ die momentane Arbeitsleistung des Systems
bestimmen durch 
\[P(t)=f(t)\cdot\dot{x}(t).\]
Passivität des Systems soll nun bedeuten, dass bei einer periodischen von
aussen anregenden Kraft, die mittlere von aussen geleistete Arbeit nicht
negativ ist.

Die Antwort auf eine periodische äußere Kraft
$f(t)=\Re\{f_0e^{-izt}\}$ mit $\Im\{z\}>0$ erhalten
wir aus (\ref{eq:Response}) zu $x(t)=\operatorname{Re}\{x_0e^{-izt}\}$, wobei
\[
  \frac{x_0}{f_0} = \chi(z) = \int_{-\infty}^\infty G(\tau)e^{iz\tau}d\tau 
\]
wie bereits in (\ref{eq:suszept}) gesehen. Das Integral ersteckt sich
weiterhin der Kausalität wegen nur auf positive Zeiten $\tau>0$. Wir schreiben
die Inverse Laplacetransformation der Suszeptibilität als
\begin{equation}
  G(\tau)=\frac{1}{2\pi}\int\limits_{-\infty+i\eta}^{+\infty+i\eta}e^{-iz\tau}\chi(z)dz=
  \frac{1}{2\pi}\int\limits_{-\infty}^{+\infty}e^{-i(\omega+i\eta)\tau}\chi(\omega+i\eta)d\omega,
  \quad \eta>0
  \label{eq:LinRespRueck}
\end{equation}
Den Beweis hierfür schauen wir uns im nächsten Abschnitt genauer an.

Zunächst benutzen wir die Zerlegung der Suszeptibilität nach Real- und
Imaginärteil, mit dem Grnzübergang, wie wir das schon im vorigen Abschnitt
getan haben
\[ 
\lim_{\eta\rightarrow 0}\chi(\omega+i\eta)=\chi'(\omega)+i\chi''(\omega)
\]
Dies erlaubt uns die Lösung $x(t)=\Re\{x_0e^{-izt}\}$ in zwei
verschiedenphasige Komponenten zu zerlegen
\begin{equation}
  x(t)=\Re\{x_0e^{-izt}\}=
  |f_0|e^{\eta t}\left(\chi'(\omega)\cos(\omega t-\varphi)
  +\chi''(\omega)sin(\omega t-\varphi)\right)
  \label{eq:Zweiphasig}
\end{equation}
Da $\chi(z)$ die Laplacetransformierte einer reellen Funktion ist, gilt
\[ [\chi(\omega+i\eta)]^*=[\chi(-\omega+i\eta)\]
Daraus wiederum folgt, dass der Realteil $\chi'$ eine gerade und der
Imaginärteil $\chi''$ eine ungerade Funktion von $\omega$ ist
\[ \chi'(-\omega)=\chi'(\omega),\quad \chi''(-\omega)=-\chi''(\omega) \]
Wir berechnen die mittlere Energieabsorption des Systems pro Zeiteinheit bei
einer periodischen Anregung mit einer äußeren Kraft der Frequenz
$\nu=\omega/2\pi$ und der Amplitude $f_0$:
\[
  \overline{P}=\frac{1}{T}\int_0^T f(t)\dot{x}(t)dt=
  \frac{1}{2}|f_0|^2\omega\chi''(\omega)
\]
Damit wird die Bedingung für die Passivität $\omega\chi''(\omega)\ge 0$.
%\section{Die Dispersionsrelationen}
%Mit Hilfe der Funktionentheorie läßt sich eine allgemeingültiger Zusammenhang
%herstellen. Es handelt sich hirbei um die \textit{Dispersionrelationen} oder
%auch \textit{Kramers-Kronig Relationen} genannt. Sie verknüpfen Real- und
%Imaginärteil der Suszeptibilität $\chi(\omega)$. Das bedeutet, dass das
%Verhalten eines linearen Systems durch eine einzige der Funktionen $G(\tau)$,
%$\chi'(\omega)$ oder $\chi''(\omega)$ vollständig bestimmt ist.
%
%Wir gehen von der dynamischen Suszeptibilität 
%\[ \chi(z) = \int_{-\infty}^\infty G(\tau)e^{iz\tau}d\tau \]
%aus. Der Integrationsbereich ist auf positive Zeiten $\tau$ beschränkt - wir
%erinnern uns, dass wir bei der Integration über $t'$ in (\ref{eq:Response})
%$t-t'=\tau$ gesetzt haben und die Funktion $G(t-t')$ nur für $t'<t$ von Null
%verschieden war, d.h. das System ist kausal - und der Expontentialfaktor drückt
%den Integranden für $\tau\rightarrow\infty$ gegen Null (wir haben $z=\omega
%+i\eta$). Damit ist 
%

\input{chapters/7_NLDGL}
\input{chapters/8_PDGL}
\chapter{Numerische Lösungsmethoden}
Nichtlineare Differentialgleichungen sind im Allgemeinen nicht analytisch
lösbar. D.h.\ es ist nicht möglich eine geschlossene Lösung anzugeben. Daher
ist es sinnvoll, sich mit numerischen Lösungsmethoden zu beschäftigen.  Wir
wollen hier nur eine begrenzte auswahl ansprechen. Für ingneieurstechnische
Anwendungen steht ganz klar das Verhältnis zwischen Aufwand für deren
Erarbeitung und Genauigkeit der Ergebnisse im Vordergrund. Deshalb gibt es
keine absolute Aussage, welche Methode die beste ist.

Wir gehen von einem Anfangswertproblem einer Differentialgleichung 1. Ordnung
und 1. Grades aus, wobei hierin auch System erster Ordnung eingeschlossen
seien. Wir schreiben das Anfangswertproblem eines Systems, wie es in
(\ref{eq:NLDGLSystem}) dargestellt ist als
\begin{align}
  \dot{\mathbf{y}}(t) =& \mathbf{f}(\mathbf{y}(t),t) \label{eq:yNLSystem}\\
  \mathbf{y}(0)       =& \mathbf{y}_{0}\nonumber
\end{align}
\section{Die Eulersche Methode}
Wir nehmen an $y(t)$ sei gekannt, dann können wir $y(t+\Delta t)$ in eine
Potenzreihe entwickeln und erhalten bei Vernachlässigung aller Terme 2. und
höherer Potenzen in $\Delta t$ 
\begin{equation}
  y(t+\Delta t) = y(t)+\Delta t f(y(t),t)+O(\Delta t^2),
  \label{eq:Euler}
\end{equation}
wobei $O(\Delta t^2)$ für quadratische und höhere Terme steht.

\begin{note}{Fehler}
  Keine Angabe eines Algorithmus zur Lösung einer DGL ist sinnvoll, wenn nicht
  gleichzeitig auch zumindest die Fehlerordnung, wenn nicht gar der lokale
  Dikretisisierungsfehler angegeben ist.
\end{note}
Wie groß ist der Fehler, den wir bei der Lösung machen, wenn wir
(\ref{eq:Euler}) für gleiche Zeitabstände zu diskreten Zeitpunkten
$t_i$ anwenden? Zumindest in einem Zeitschritt gelingt uns sofort eine Aussage.
Angenommen $Y_i$ sei die exakte Lösung der Differentialgleichung zum Zeitpunkt
$t_i$. Dann können wir $Y_{i+1}$ ebenfalls durch eine Taylorreihe darstellen
\begin{equation}
  Y_{i+1}=Y_{i}
  +\underbrace{\Delta t \left(\frac{dY}{dt}\right)_i}_{f(Y_i,t_i)}
  +\underbrace{\frac{\Delta t^2}{2}\left(\frac{d^2Y}{dt^2}\right)_i}_{\dot{f}(Y_i,t_i)}+\dots
  \label{eq:TaylorExact}
\end{equation}
Wobei wir jetzt auf der rechten Seite von (\ref{eq:TaylorExact}) $Y_i$ durch
$y_i$ ersetzen und von (\ref{eq:Euler}) die Potenzreihe (\ref{eq:TaylorExact})
abziehen. Damit erhalten wir für den Fehler
\begin{equation}
  y_{i+1}-Y_{i+1}=E_{i+1}=-\frac{\Delta t^2}{2}\dot{f}(Y_i,t_i)O(\Delta t^3)
  \label{eq:LokalerFehler}
\end{equation}
Dies ist der Fehler pro Zeitschritt, auch lokaler Diskretisierungfehler genannt.
\section{Runge Kutta Methode}
Wir gehen von einer zentralen finiten Differenzenformel aus. Das bedeutet, dass
die finite Differenz die Ableitung zum Zeitpunkt $t_i+\frac{\Delta t}{2}$
nähert. Wir erhalten
\begin{equation}
  \frac{y_{i+1}- y_i}{\Delta}=\dot{y}_{i+1/2}= f(y_{i+1/2},t_{i+1/2})
  \label{eq:Midpoint}
\end{equation}
wobei aber $y_{i+1/2}$ unbekannt ist. Dies beschaffen wir uns durch
Vorwärtsintegration um einen halben Zeitschritt $\frac{\Delta t}{2}$, somit
\begin{align*}
  \hat{y}_{i+1/2}=&y_i+\frac{\Delta t}{2}f(y_{i},t_{i})\\
  t_{i+1/2} =& t_i+\frac{\Delta t}{2}\\
  y_{i+1} =& y_{i}+\Delta t f(\hat{y}_{i+1/2},t_{i+1/2})
\end{align*}
Dies ist das sogenannte Runge-Kutta-Schema 2. Ordnung. Wir sehen, dass wir
hier, im Gegensatz zum Euler-Methode die Funktion $f$ zweimal auswerten müssen.

Ein verbessertes Schema ist die Runge-Kutta-Methode 4. Ordnung. Wie der name
sagt, müssen wir hier die  Funktion $f$ viermal auswerten. Wir erhalten
\begin{align}\label{eq:RK4order}
  y_{i+1} =& y_{i}+\frac{\Delta t}{6}\left( 
    f(y_i,t_i)+2f(\hat{y}_{i+1/2},t_{i+1/2})+2f(\hat{\hat{y}}_{i+1/2},t_{i+1/2})
  +f(\hat{y}_{i+1},t_{i+1})\right)\\
  \hat{y}_{i+1/2}=&y_i+\frac{\Delta t}{2}f(y_{i},t_{i})\nonumber\\
  \hat{\hat{y}}_{i+1/2}=&y_i+\frac{\Delta t}{2}f(y_{i+1/2},t_{i+1/2})\nonumber\\
  \hat{y}_{i+1}=&y_i+\Delta tf(\hat{\hat{y}}_{i+1/2},t_{i+1/2})\nonumber
\end{align}

\section{Mehrschrittmethoden}
Die Idee hinter den Mehrschrittmethoden ist, (\ref{eq:NLDGLSystem}) zwischen
$t$ und $t+\Delta t$ zu integrieren
\begin{equation*}
y(t_{n+s})=y(t_{n+s-1})+\int\limits_{t_{n+s}}^{t_{n+s-1}}f(y(t'),t')dt'
\end{equation*}
Das Integral gilt es nun zu nähern. Wir haben ja bereits Werte für $y(t_i)$ zu
den Zeitpunkten $t_n$  bis $t_{n+s-1}$ berechnet. Wir nehmen  diese um eine
Interpolationsfunktion aufzustellen, die wir nun in den Grenzen von $t_{n+s-1}$
bis $t_{n+s}$ integrieren können. Wir nähern $f(y(t),t)$ z.B. durch ein
Polynom $p(t)$ der Ordnung $s-1$, das wir dergestalt konstruiert haben, dass es an
den Interpolationspunkten mit $f(y(t),t)$ übereinstimmt

Das Interpolationspolynom ist analytisch integrierbar.  Explizit schreiben wir
das folgendermaßen
\[ p(t)=\sum\limits_{m=0}^{s-1}p_m(t)f(y(t_{n+m},t_{n+m}). \]
Hierbei bezeichne $p_m(t)$ ein Polynom, das zum Zeitpunkt $t_{n+m}$ gleich 1
ist und zu allen anderen Zeitpunkten 0. Damit nutzen wir $s$ vorhergehende
Zeitschritte aus. Daher auch der Name Mehrschrittverfahren und hier im
speziellen beschrieben ist die Methode von Adams.
%\subsection{Predictor-Corrector Methode}
%\subsection{Shooting- und Matrixmethode für Randwertprobleme}
%\input{chapters/10_Modellidentifikation}
\input{chapters/11_Variationsrechnung}

%
\begin{appendices}
\noappendicestocpagenum
\renewcommand\appendixtocname{Anhang}
\renewcommand\appendixpagename{Anhang}
\input{chapters/A_Residuensatz}
\input{chapters/B_VNOS}
\chapter{Interpolation und Kurvenanpassung}

\section{Least square fit}
Suppose you are given a set of data $z_i\in\{Z\}$, e.g. the $i=1,\dots,29$ indicating the channels of an EEG measurement, measured at points $x_i$ and $y_i$. We want to find a function $f(x,y;\mathbf{p})$ which depends on a vector $\mathbf{p})$ of paramters $p_k$. 
\subsection{Polynomial basis}
Assume that $f(x=0,y=0)=0$. Thus we propose a polynomial function $f(x,y;\mathbf{p})=p_0\cdot x+p_1\cdot y+p_2\cdot x^2+p_3\cdot x\cdot y+p_4\cdot y^2$ the data set of measurements $\{x_i,y_i,z_i\}$ we want to fit to. We choose a least square method, i.e.
\[
\sum_i \left( f(x_i,y_i;\mathbf{p})-z_i\right)^2 = \text{Min}.
\]
Which in turn is equivalent to require
\[
\frac{\partial}{\partial p_k}\sum_i \left( f(x_i,y_i;\mathbf{p})-z_i\right)^2=\sum_i \left( f(x_i,y_i;\mathbf{p})-z_i\right)\frac{\partial}{\partial p_k}f(x_i,y_i;\mathbf{p})=0
\]
this yields a linear system of equation for the $p_k$ to read
\begin{align*}
p_0\sum_i x_i^2+p_1\sum_i x_i\cdot y_i+p_2\sum_i x_i^3+p_3\sum_i x_i^2\cdot y_i+p_4\sum_i x_i\cdot y_i^2&=\sum_iz_i\cdot x_i\\
p_0\sum_i x_i\cdot y_i+p_1\sum_i y_i^2+p_2\sum_i x_i^2\cdot y_i+p_3\sum_i x_i\cdot y_i^2+p_4\sum_i y_i^3&=\sum_iz_i\cdot y_i\\
p_0\sum_i x_i^3+p_1\sum_i x_i^2\cdot y_i+p_2\sum_i x_i^4+p_3\sum_i x_i^3\cdot y_i+p_4\sum_i x_i^2\cdot y_i^2&=\sum_iz_i\cdot x_i^2\\
p_0\sum_i x_i^2\cdot y_i+p_1\sum_i x_i\cdot y_i^2+p_2\sum_i x_i^3\cdot y_i+p_3\sum_i x_i^2\cdot y_i^2+p_4\sum_i x_i\cdot y_i^3&=\sum_iz_i\cdot x_i\cdot y_i\\
p_0\sum_i x_i\cdot y_i^2+p_1\sum_i y_i^3+p_2\sum_i x_i^2\cdot y_i^2+p_3\sum_i x_i\cdot y_i^3+p_4\sum_i x_i\cdot y_i^4&=\sum_iz_i\cdot y_i^2
\end{align*}
Let us write this as a 
\[
\mathbf{A}\cdot\mathbf{p}=\mathbf{b}
\]
Note that $\mathbf{A}$ is symmetric and the matrix entries read
\[ 
\mathbf{A}=
\begin{pmatrix}
\sum_i x_i^2&\sum_i x_i\cdot y_i&\sum_i x_i^3&\sum_i x_i^2\cdot y_i&\sum_i x_i\cdot y_i^2\\
\sum_i x_i\cdot y_i&\sum_i y_i^2&\sum_i x_i^2\cdot y_i&\sum_i x_i\cdot y_i^2&\sum_i y_i^3\\
\sum_i x_i^3&\sum_i x_i^2\cdot y_i&\sum_i x_i^4&\sum_i x_i^3\cdot y_i&\sum_i x_i^2\cdot y_i^2\\
\sum_i x_i^2\cdot y_i&\sum_i x_i\cdot y_i^2&\sum_i x_i^3\cdot y_i&\sum_i x_i^2\cdot y_i^2&\sum_i x_i\cdot y_i^3\\
\sum_i x_i\cdot y_i^2&\sum_i y_i^3&\sum_i x_i^2\cdot y_i^2&\sum_i x_i\cdot y_i^3&\sum_i x_i\cdot y_i^4
\end{pmatrix},
\]
while we habe for the r.h.s. of the linear system
\[ 
\mathbf{b}=
\begin{pmatrix}
\sum_iz_i\cdot x_i\\
\sum_iz_i\cdot y_i\\
\sum_iz_i\cdot x_i^2\\
\sum_iz_i\cdot x_i\cdot y_i\\
\sum_iz_i\cdot y_i^2
\end{pmatrix},
\]
$\mathbf{A}$ and $\mathbf{b}$ can be easily calculated and thus we have $\mathbf{p}=\mathbf{A}^{-1}\cdot\mathbf{b}$. The only task left is to find an analytical expression for the inverse of $\mathbf{A}$, which components read
\[
\mathbf{A}^{-1}_{ij}=\frac{\alpha_{ij}}{|\mathbf{A}|},
\]
where the $\alpha_{ij}$ are the elements of the adjoint matrix $\mathbf{A}_{ad}$. 
\begin{note}{The adjoint matrix elements}
\[
\alpha_{ij}=(-1)^{i+j}
\left|
\begin{matrix}
a_{11}&\dots &a_{1,j-1}& &a_{1,j+1}&\dots &a_{1n}\\
a_{11}&\dots &a_{1,j-1}& &a_{1,j+1}&\dots &a_{1n}\\
&&&\vdots &&&\\
a_{11}&\dots &a_{1,j-1}& &a_{1,j+1}&\dots &a_{1n}\\
a_{11}&\dots &a_{1,j-1}& &a_{1,j+1}&\dots &a_{1n}\\
\end{matrix}
\right|\text{, where }a_{ij}=(\mathbf{A})_{ij}
\]
\end{note}
\subsection{Spherical harmonics function as a basis}
Now let $z_i\in\{Z\}$ be the $i=1,\dots,29$ indicating the channels of an EEG measurement at positions $\theta_i(t),\varphi_i(t)$, i.e. 29 time series'.
We propose a set of spherical harmonics functions $f(x,y;\mathbf{p})=Y_{nl}(\theta,\varphi)$ the data set of measurements $\{\theta_i,\varphi_i,z_i\}$ we want to fit to. The least square method states that
\[
\sum_i \left(\sum\limits_{nl}c_{nl}Y_{nl}(\theta_i,\varphi_i)-z_i\right)^2 = \text{Min}.
\]
Thus all the derivatives w.r.t. $c_{mk}$ must vanish, to read
\begin{equation}\label{eq:lstsqsph}
    \sum\limits_i2\left(\sum\limits_{nl}c_{nl}Y_{nl}(\theta_i,\varphi_i)-z_i\right)\cdot Y_{mk}(\theta_i,\varphi_i)=0
\end{equation}
For the spherical harmonics we have $l=-n,\dots,n$ for every $n$.
\begin{example}{Spherical harmonics least square fit}
Let us limit the number of spherical harmonics in (\ref{eq:lstsqsph}) to $n=0,1$. Thus we are left with 4 parameters $c_0=c_0,c_{1,-1}=c_1,c_{1,0}=c_2,c_{1,1}=c_4$. Thus (\ref{eq:lstsqsph}) reads
\[
\sum\limits_{i=0}^{29}
\left(\sum\limits_{j=0}c_{j}Y_{j}(\theta_i,\varphi_i)
-z_i\right)\cdot Y_{k}(\theta_i,\varphi_i)=0
\]
We get four equations - remember $k=0,1,2,3$
\begin{align*}
&c_{0}\sum\limits_{i}Y_{0}(\theta_i,\varphi_i)\cdot Y_{k}(\theta_i,\varphi_i)+
c_{1}\sum\limits_{i}Y_{1}(\theta_i,\varphi_i)\cdot Y_{k}(\theta_i,\varphi_i)+\\
&c_{2}\sum\limits_{i}Y_{2}(\theta_i,\varphi_i)\cdot Y_{k}(\theta_i,\varphi_i)+
c_{3}\sum\limits_{i}Y_{3}(\theta_i,\varphi_i)\cdot Y_{k}(\theta_i,\varphi_i)
=\sum\limits_{i}z_iY_{k}(\theta_i,\varphi_i)
\end{align*}
The $a_{kj}=\sum\limits_{i}Y_{j}(\theta_i,\varphi_i)\cdot Y_{k}(\theta_i,\varphi_i)$ and the $b_k=\sum\limits_{i}z_iY_{k}(\theta_i,\varphi_i)$ are just real numbers. Thus we end up with the linear system
\begin{align*}
    a_{00}c_0+a_{01}c_1+a_{02}c_2+a_{03}c_3=&b_0\\
    a_{10}c_0+a_{11}c_1+a_{12}c_2+a_{13}c_3=&b_1\\
    a_{20}c_0+a_{21}c_1+a_{22}c_2+a_{23}c_3=&b_2\\
    a_{30}c_0+a_{31}c_1+a_{32}c_2+a_{33}c_3=&b_3\\
\end{align*}
to determine the $c_i$.
\end{example}
\end{appendices}
\end{document}